\documentclass[aip,preprint]{revtex4-1}

\usepackage[utf8]{inputenc}
\usepackage{amsmath}
\usepackage{amssymb}
\usepackage{braket}
\usepackage{xspace}

\usepackage[capitalize]{cleveref}

\newcommand{\dd}{\mathrm{d}}
\newcommand{\ao}{\mathrm{AO}}
\newcommand{\aux}{\mathrm{aux}}
\newcommand{\core}{\mathrm{core}}
\newcommand{\act}{\mathrm{act}}
\newcommand{\sa}{\mathrm{SA}}
\newcommand{\df}{\mathrm{DF}}
\newcommand{\hf}{\mathrm{HF}}
\newcommand{\pert}{x}

\begin{document}

\title{Technical Notes: Gradient Theory}

\author{Forte2 Developers}

\date{\today}

\maketitle

\section{Scope}
This note covers the technical details of implementing analytic nuclear gradients, specifically in the case where electron-repulsion integrals are represented using the density-fitting approximation. 

\section{Notation}
Let $\chi_\mu(\mathbf{r})$ denote an orbital basis function and $\chi_P(\mathbf{r})$ denote an auxiliary
basis function.
The one-electron integrals in the orbital basis are given by:
\begin{equation}
h_{\mu \nu}
= \braket{\chi_\mu|\hat{h}|\chi_\nu}
= \int
\chi_\mu(\mathbf{r})
\hat{h}
\chi_\nu(\mathbf{r})\dd\mathbf{r},
\end{equation}
where the one-body operator is the sum of the kinetic and electron-nuclear potential:
\begin{equation}
\hat{h} = - \frac{1}{2} \nabla^2 + V_\mathrm{eN}(\mathbf{r})
= - \frac{1}{2} \nabla^2 - \sum_{A} \frac{Z_A}{|\mathbf{R}_A - \mathbf{r}|}.
\end{equation}
The two-electron atomic integrals in chemists' notation are defined as:
\begin{equation}
(\mu\nu|\rho\sigma)
= \iint
\chi_\mu(\mathbf{r}_1)\chi_\nu(\mathbf{r}_1)
\frac{1}{r_{12}}
\chi_\rho(\mathbf{r}_2)\chi_\sigma(\mathbf{r}_2)
\dd\mathbf{r}_1\dd\mathbf{r}_2,
\end{equation}
and we also define the three-center tensor and Coulomb metric as
\begin{align}
J^P_{\mu\nu} &\equiv (P|\mu\nu) = \iint
\chi_P^*(\mathbf{r}_1)\chi_\nu(\mathbf{r}_1)
\frac{1}{r_{12}}
\chi_\rho(\mathbf{r}_2)^*\chi_\sigma(\mathbf{r}_2)
\dd\mathbf{r}_1\dd\mathbf{r}_2 , \\
M_{PQ} &\equiv (P|Q) = \iint
\chi_P^*(\mathbf{r}_1)
\frac{1}{r_{12}}
\chi_Q(\mathbf{r}_2)
\dd\mathbf{r}_1\dd\mathbf{r}_2.
\end{align}
The inverse Coulomb metric is denoted as:
\begin{equation}
N_{PQ} \equiv (M^{-1})_{PQ}.
\end{equation}
When writing summations, the range will be implicitly determined by the type of index.
Following this convention, a density-fitted four-center integral is written as:
\begin{equation}
(\mu\nu|\rho\sigma)_\df
= \sum_{PQ} J^P_{\mu\nu} N_{PQ} J^Q_{\rho\sigma}
= \sum_{Q} b^{P}_{\mu \nu} b^{P}_{\rho\sigma},
\label{eq:df_eri}
\end{equation}
where the three-index tensor $b^{P}_{\mu \nu}$ is defined as:
\begin{equation}
b^{P}_{\mu \nu} = \sum_{Q} J^Q_{\mu\nu} [\mathbf{N}^{1/2}]_{QP}.
\end{equation}

Let $x$ denote one Cartesian nuclear coordinate, for example
$x = R_{A\alpha}$ with atom $A$ and component
$\alpha \in \{x,y,z\}$.  A superscript $\pert$ denotes differentiation with
respect to $x$:
\begin{equation}
A^{\pert} \equiv \frac{\partial A}{\partial x}.
\end{equation}
The one-electron operator depends explicitly on the nuclear coordinates, with $h^{\pert}_{\mu\nu}$ being expressed as
\begin{equation}
h^{\pert}_{\mu\nu} = 
\braket{\chi^{\pert}_\mu|\hat{h}|\chi_\nu}
+ \braket{\chi_\mu|\hat{h}^{\pert}|\chi_\nu} 
+ \braket{\chi_\mu|\hat{h}|\chi^{\pert}_\nu}.
\end{equation}
The electron-electron operator has no explicit dependence on nuclear
coordinates in the clamped-nuclei electronic Hamiltonian.  Therefore
$(\mu\nu|\rho\sigma)^{\pert}$, $J^{P,\pert}_{\mu\nu}$, and $M^{\pert}_{PQ}$ arise from
the nuclear dependence of the basis functions and auxiliary basis functions.

\section{Hartree--Fock Gradient Theory}

To derive gradient expressions for the Hartree--Fock method, we begin by defining a Lagrangian that includes the energy, the diagonality of the Fock matrix, and the orthonormality of the occupied orbitals:
\begin{equation}
\begin{split}
\mathcal{L}_\hf(\mathbf{R},\mathbf{C},\mathbf{C}^*,\boldsymbol{\epsilon},\mathbf{z},,\boldsymbol{\omega})
= & E_\mathrm{HF}(\mathbf{R},\mathbf{C},\mathbf{C}^*) \\
& + \sum_{ij}^{\mathbb{O}} z_{ji} (f_{ij}(\mathbf{R},\mathbf{C},\mathbf{C}^*) - \delta_{ij}\epsilon_i(\mathbf{R})) \\
& - \sum_{ij}^{\mathbb{O}} \omega_{ji} (S_{ij}(\mathbf{R},\mathbf{C},\mathbf{C}^*) - \delta_{ij}),
\end{split}
\end{equation}
where
\begin{equation}
E_\hf(\mathbf{R},\mathbf{C},\mathbf{C}^*)
= V_\mathrm{NN}
+ \sum_{i}^{\mathbb{O}} h_{ii}
+ \frac{1}{2} \sum_{ij}^{\mathbb{O}} \braket{ij\|ij},
\end{equation}
\begin{equation}
f_{ij}(\mathbf{R},\mathbf{C},\mathbf{C}^*) = h_{ij} +\sum_{k}^{\mathbb{O}} \braket{ik\|jk},
\end{equation}
and
\begin{equation}
S_{ij}(\mathbf{R},\mathbf{C},\mathbf{C}^*) = \braket{\phi_i(\mathbf{R},\mathbf{C},\mathbf{C}^*)|\phi_j(\mathbf{R},\mathbf{C},\mathbf{C}^*)}.
\end{equation}

In defining the Lagrangian we consider the nuclear coordinates ($\mathbf{R}$), the orbital coefficients ($\mathbf{C},\mathbf{C}^*$), the orbital energies ($\boldsymbol{\epsilon}$), the orbital rotation parameters ($\mathbf{z}$), and  the orthogonality parameters ($\boldsymbol{\omega}$) to be \textbf{independent parameters}.
All other quantities ($E_\hf$, $f_{ij}$, \ldots) are dependent on these parameters.
Lastly, the variable ($x$) indicates an external perturbation.

We assume the most general case, and consider complex orbital coefficients and Fock matrices and overlap integral matrices that are complex and Hermitian.
To ensure that $\mathcal{L}_\hf$ is a real function (otherwise we cannot define its stationary point!), the matrices $\mathbf{z}$ and  $\boldsymbol{\omega}$ satisfy are Hermitian, that is $\mathbf{z}^\dagger = \mathbf{z}$ and $\boldsymbol{\omega}^\dagger = \boldsymbol{\omega}$.
To evaluate derivatives with respect to the orbital coefficients, we use Wirtinger calculus, which treats the pair of matrices $\mathbf{C}$ and its complex conjugate $\mathbf{C}^*$ as independent variables.

Next, we impose all the constraints compatible with the Hartree--Fock equations by imposing stationary conditions for all independent parameters.
Note that stationarity with respect to some parameter $\lambda$ is imposed holding all other parameters fixed. We do not indicate this explicitly in the equations below.
We will also omit the parameter dependence of all quantities.
Stationarity of $\mathcal{L}_\mathrm{HF}$ with respect to the orbital rotation parameters and orthogonality parameters just gives us back the corresponding constraints:
\begin{align}
\frac{\partial \mathcal{L}_\hf}{\partial z_{ji}} = 0 \quad & \Rightarrow \quad f_{ij} = \delta_{ij}\epsilon_i, \\
\frac{\partial \mathcal{L}_\hf}{\partial \omega_{ji}} = 0 \quad &\Rightarrow \quad \braket{\phi_i|\phi_j} = \delta_{ij}.
\end{align}

Nontrivial stationarity conditions start to appear when we make $\mathcal{L}_\hf$ stationary with respect to the orbital energies:
\begin{equation}
\frac{\partial \mathcal{L}_\hf}{\partial \epsilon_{i}} = 0 \quad \Rightarrow \quad z_{ii} = 0.
\end{equation}
To impose stationarity with respect to the orbital coefficients, we introduce a directional derivative (keeping all other independent variables constant):
\begin{equation}
D_{ij} = \sum_{\mu} C_{\mu i}\frac{\partial}{\partial C_{\mu j}},
\end{equation}
which we use to impose stationarity conditions in the following way:
\begin{equation}
X_{ij} = D_{ij}\mathcal{L}_\hf = \sum_{\mu} C_{\mu i}\frac{\partial \mathcal{L}_\hf}{\partial C_{\mu j}} = 0 \quad \forall i,j \in \mathbb{O}.
\end{equation}
The directional derivative acts only on kets, replacing a spin orbital and introducing a Kronecker delta:
\begin{equation}
D_{ij} \ket{\psi_k} = 
\sum_{\mu} C_{\mu i}\frac{\partial \ket{\psi_k}}{\partial C_{\mu j}}
= \sum_{\mu} C_{\mu i}\ket{\chi_\mu}\delta_{jk}
= \ket{\psi_i}\delta_{jk}.
\end{equation}
When $D_{ij}$ is applied to the Fock matrix it gives:
\begin{equation}
D_{ij} 	f_{kl}
= D_{ij} \braket{\psi_k| \hat{h} | \psi_l} +\sum_{m}^{\mathbb{O}} D_{ij} \braket{km\|lm}
= \underbrace{\delta_{jl} h_{ki}
+\sum_{m}^{\mathbb{O}} \braket{km\|im} \delta_{jl}}_{\delta_{jl} f_{ki}}
+\braket{kj\|li}.
\end{equation}
Next we consider the contribution to the energy, starting with the one-body term:
\begin{equation}
D_{ij}
(\sum_{k}^{\mathbb{O}} h_{kk})
=
\sum_{k}^{\mathbb{O}} D_{ij}\braket{\psi_k| \hat{h} | \psi_k}
=
\braket{\psi_j| \hat{h} | \psi_i},
\end{equation}
and the two-body term:
\begin{equation}
D_{ij}
(\frac{1}{2}\sum_{kl}^{\mathbb{O}} \braket{kl\|kl})
= \sum_{k}^{\mathbb{O}} \braket{jk\|ik},
\end{equation}
which yields
\begin{equation}
D_{ij}
E_\hf
= h_{ji} +\sum_{k}^{\mathbb{O}} \braket{jk\|ik} = f_{ji} = \delta_{ij} \epsilon_{i}.
\end{equation}
The derivative of the diagonality constraint is:
\begin{equation}
D_{ij}
\sum_{kl}^{\mathbb{O}} z_{lk} (f_{kl} - \delta_{kl}\epsilon_k)
= \sum_{kl}^{\mathbb{O}} z_{lk} D_{ij} f_{kl} = 
z_{ji} \epsilon_{i}
+\sum_{kl}^{\mathbb{O}} z_{lk}\braket{kj\|li}.
\end{equation}
The derivative of the orthonormality constraint is:
\begin{equation}
D_{ij}
\sum_{kl}^{\mathbb{O}} \omega_{lk} (\braket{\phi_k | \phi_l} - \delta_{kl})
=
\sum_{kl}^{\mathbb{O}} \omega_{lk} D_{ij}\braket{\phi_k | \phi_l}
=
\sum_{kl}^{\mathbb{O}} \omega_{lk}\braket{\phi_k | \phi_i} \delta_{jl}
= \omega_{ji}.
\end{equation}

Putting all together:
\begin{equation}
X_{ij} =\delta_{ij} \epsilon_{i} + z_{ji} \epsilon_{i}
+\sum_{kl}^{\mathbb{O}} z_{lk}\braket{kj\|li}
- \omega_{ji} = 0.
\end{equation}
Using the fact that $z$ and $\omega$ are Hermitian we can evaluate the difference:
\begin{equation}
X_{ij} - X_{ji}^{*} = z_{ji} \epsilon_{i} -  z_{ij}^* \epsilon_{j} = 0, \Rightarrow z_{ji}(\epsilon_{i} - \epsilon_{j})=0.
\end{equation}
When $\epsilon_{i} - \epsilon_{j} \neq 0$, we find the solution:
\begin{equation}
z_{ji} = 0.	
\end{equation}
When $\epsilon_{i} = \epsilon_{j}$, we instead fall into the case of a rotation among two degenerate orbitals. In this case, any rotation keeps the Fock matrix diagonal, so we have a gauge degree of freedom. A simple choice is to set for all orbital rotations:
\begin{equation}
z_{ji} = 0.
\end{equation}
Substituting in the equation for $X_{ij}$ we arrive at:
\begin{equation}
X_{ij} =\delta_{ij} \epsilon_{i}
- \omega_{ji} = 0, \Rightarrow
\omega_{ji} = \delta_{ij} \epsilon_{i}.
\end{equation}

Our next step should be considering the stationarity with respect to the complex conjugate of the orbital coefficients $\textbf{C}^*$.
However, for a real Lagrangian, the stationarity equations for $\textbf{C}^*$ are the complex conjugates of the stationarity equations for $\textbf{C}$. Therefore, they contain no independent information, and we can skip the derivation.

Finally, to obtain the derivative, we evaluate the derivative of the Lagrangian with respect to $\pert$ while holding all other parameters fixed.
We then arrive at
\begin{equation}
\mathcal{L}_\hf^{\pert}
= E^{\pert}_\mathrm{HF}
+ \sum_{ij}^{\mathbb{O}} z_{ji} f_{ij}^{\pert}
- \sum_{ij}^{\mathbb{O}} \omega_{ji} \braket{\phi_i|\phi_j}^{\pert},
\end{equation}
where $x$ stands for the derivative with respect to $x$ with all independent parameters kept constant.
Therefore, the expressions involving $x$ will only contain integral derivatives, and no parameter derivatives.
Inserting our solutions for the Lagrange multipliers ($\mathbf{z} = 0$, $\omega_{ji} = \delta_{ij} \epsilon_{i}$), we arrive at our final gradient expression:
\begin{equation}
\mathcal{L}_\hf^{\pert}
= E^{\pert}_\mathrm{HF}
- \sum_{i}^{\mathbb{O}} \epsilon_{i} \braket{\phi_i|\phi_i}^{\pert}.
\end{equation}

\section{Energy Expectation Values and Density Matrices}

We now proceed to develop the general theory for evaluating gradients using density fitting.

\subsection{Spinorbital Form}

Let the indices $p,q,r,s$ denote general spinorbitals.  In a spinorbital basis, write the electronic
Hamiltonian as:
\begin{equation}
\hat{H}
=
V_\mathrm{NN} +
\sum_{pq}^{\mathbb{G}} h_{pq}\hat{a}^\dagger_p\hat{a}_q
+
\frac{1}{4}
\sum_{pqrs}^{\mathbb{G}} \braket{pq\|rs}
\hat{a}^\dagger_p\hat{a}^\dagger_q\hat{a}_s\hat{a}_r.
\label{eq:spinorbital_hamiltonian}
\end{equation}
Here we use antisymmetrized integrals defined as:
\begin{equation}
\braket{pq\|rs} = \braket{pq|rs} - \braket{pq|sr},
\end{equation}
where the integral $\braket{pq|rs}$ assumes the following convention:
\begin{equation}
\braket{pq|rs}
= \iint
\psi^*_p(x_1) \psi^*_q(x_2)
\frac{1}{r_{12}}
\psi_r(x_1) \psi_s(x_2)
\dd x_1\dd x_2.
\end{equation}
We define the spinorbital one- and two-particle density matrices as
\begin{align}
\gamma_{pq}
&=
\langle \Psi |
\hat{a}^\dagger_p\hat{a}_q
| \Psi \rangle,
\label{eq:spinorbital_1rdm}\\
\gamma_{pq,rs}
&=
\langle \Psi |
\hat{a}^\dagger_p\hat{a}^\dagger_q\hat{a}_s\hat{a}_r
| \Psi \rangle.
\label{eq:spinorbital_2rdm}
\end{align}
Then
\begin{equation}
E
= V_\mathrm{NN} +
\sum_{pq}^{\mathbb{G}} h_{pq}\gamma_{pq}
+
\frac{1}{4}
\sum_{pqrs}^{\mathbb{G}} \braket{pq\|rs}\gamma_{pq,rs}.
\label{eq:spinorbital_energy}
\end{equation}

For \textbf{a single Slater determinant} the 2-RDM factorizes into a product of two 1-RDMs:
\begin{equation}
\gamma_{pq,rs}
=
\gamma_{pr}\gamma_{qs}
-
\gamma_{ps}\gamma_{qr}.
\label{eq:determinant_spinorbital_2rdm}
\end{equation}
This formula is the most compact way to derive the RHF, UHF, ROHF, and GHF
pair densities used below.


\subsection{Spin-Free Spatial-Orbital Form}

For spin-free nonrelativistic methods, let $p,q,r,s$ denote spatial orbitals and
introduce the spin-free excitation operator
\begin{equation}
\hat{E}_{pq}
=
\sum_{\sigma\in\{\alpha,\beta\}}
\hat{a}^\dagger_{p\sigma}\hat{a}_{q\sigma}.
\end{equation}
The Hamiltonian can be written as
\begin{equation}
\hat{H}
= V_\mathrm{NN} +
\sum_{pq}^{\mathbb{G}}h_{pq}\hat{E}_{pq}
+
\frac{1}{2}
\sum_{pqrs}^{\mathbb{G}}\braket{pq|rs}
\left(
\hat{E}_{pr}\hat{E}_{qs}
-
\delta_{qr}\hat{E}_{ps}
\right).
\label{eq:spinfree_hamiltonian}
\end{equation}
Define the spin-free one-particle density matrix:
\begin{equation}
\Gamma_{pq}
=
\langle \Psi | \hat{E}_{pq} | \Psi \rangle
= \gamma_{p\alpha \, q\alpha} + \gamma_{p\beta \, q\beta},
\label{eq:spinfree_1rdm}
\end{equation}
and the spin-free two-particle density matrix:
\begin{equation}
\Gamma_{pq,rs}
=
\braket{ \Psi |
\hat{E}_{pr}\hat{E}_{qs}
-
\delta_{qr}\hat{E}_{ps}
| \Psi }
=
\gamma_{p\alpha \, q\alpha, r\alpha \, s\alpha }
+ \gamma_{p\alpha \, q\beta, r\alpha \, s\beta }
+ \gamma_{q\alpha \, p\beta, s\alpha \, r\beta }
+ \gamma_{p\beta \, q\beta, r\beta \, s\beta }.
\label{eq:spinfree_2rdm}
\end{equation}
The energy expectation value is
\begin{equation}
E
= V_\mathrm{NN} +
\sum_{pq}^{\mathbb{G}} h_{pq}\Gamma_{pq}
+
\frac{1}{2}
\sum_{pqrs}^{\mathbb{G}} \braket{pq|rs}\Gamma_{pq,rs}.
\label{eq:spinfree_energy}
\end{equation}

%For a determinant expressed in a spatial AO basis, define spin-block density
%matrices
%\begin{equation}
%\gamma^{\omega\tau}_{\mu\nu}
%=
%\left\langle
%\hat{a}^\dagger_{\mu\omega}\hat{a}_{\nu\tau}
%\right\rangle,
%\qquad
%\omega,\tau\in\{\alpha,\beta\},
%\label{eq:spin_block_density}
%\end{equation}
%and the charge density
%\begin{equation}
%P_{\mu\nu}
%=
%\gamma^{\alpha\alpha}_{\mu\nu}
%+
%\gamma^{\beta\beta}_{\mu\nu}.
%\label{eq:charge_density}
%\end{equation}
%Substituting the determinant spinorbital density
%\cref{eq:determinant_spinorbital_2rdm} into the spin-summed spatial integral
%contraction gives the determinant spin-free pair density
%\begin{equation}
%\Pi_{\mu\nu,\rho\sigma}
%=
%P_{\mu\nu}P_{\rho\sigma}
%-
%\sum_{\omega\tau}
%\gamma^{\omega\tau}_{\mu\sigma}
%\gamma^{\tau\omega}_{\rho\nu}.
%\label{eq:determinant_spinfree_pair_density}
%\end{equation}
%For collinear spin methods, the off-diagonal spin blocks vanish and this
%reduces to
%\begin{equation}
%\Pi_{\mu\nu,\rho\sigma}
%=
%P_{\mu\nu}P_{\rho\sigma}
%-
%\gamma^{\alpha\alpha}_{\mu\sigma}
%\gamma^{\alpha\alpha}_{\rho\nu}
%-
%\gamma^{\beta\beta}_{\mu\sigma}
%\gamma^{\beta\beta}_{\rho\nu}.
%\label{eq:collinear_pair_density}
%\end{equation}
%For RHF, $\gamma^{\alpha\alpha}=\gamma^{\beta\beta}=P/2$, so
%\begin{equation}
%\Pi_{\mu\nu,\rho\sigma}^{\mathrm{RHF}}
%=
%P_{\mu\nu}P_{\rho\sigma}
%-
%\frac{1}{2}P_{\mu\sigma}P_{\rho\nu}.
%\label{eq:rhf_pair_density}
%\end{equation}




Expanding this expression in terms of orbital basis functions we can write it as:
\begin{equation}
\begin{split}
E
& = V_\mathrm{NN} +
\sum_{\mu \nu} h_{\mu\nu} \sum_{pq}^{\mathbb{G}} C_{\mu p}^* \Gamma_{pq} C_{\nu q}
+
\frac{1}{2}
\sum_{\mu \nu \rho \sigma}
\sum_{pqrs}^{\mathbb{G}} \braket{\mu \nu |\rho \sigma }
C_{\mu p}^* C_{\nu q}^* \Gamma_{pq,rs}
C_{\rho r} C_{\sigma s} \\
& = V_\mathrm{NN} + \sum_{\mu \nu} h_{\mu\nu} \Pi_{\mu \nu}
+
\frac{1}{2}
\sum_{\mu \nu \rho \sigma}
(\mu \nu |\rho \sigma)
\Pi_{\mu \nu,\rho \sigma},
\label{eq:ao_spinorbital_energy}
\end{split}
\end{equation}
where the one- and two-electron AO densities are defined as:
\begin{equation}
\Pi_{\mu \nu} = \sum_{pq}^{\mathbb{G}} C_{\mu p}^* \Gamma_{pq} C_{\nu q},
\end{equation}
and
\begin{equation}
\Pi_{\mu \nu,\rho \sigma} = \sum_{pqrs}^{\mathbb{G}} 
C_{\mu p}^* C_{\nu q}^* \Gamma_{pr,qs}
C_{\rho r} C_{\sigma s}.
\end{equation}
\textbf{Note that in defining $\Pi_{\mu \nu,\rho \sigma}$ we had to permute the second and third index in the summation over spinorbital labels.}

\section{General Energy Expression under Density fitting}

Under the density-fitted approximation we can write:
\begin{equation}
E = V_\mathrm{NN} +
\sum_{\mu \nu} h_{\mu\nu} \Pi_{\mu \nu}
+
\frac{1}{2}
\sum_{\mu \nu \rho \sigma}
(\mu \nu |\rho \sigma)_\df
\Pi_{\mu \nu,\rho \sigma}.
\label{eq:e2_general}
\end{equation}
%For a variational stationary wave function, $\Pi$ is the two-particle density appearing in the energy Lagrangian.  For state-averaged CASSCF, it is the state-weighted average pair density for the state-averaged objective.  For a nonstationary target state evaluated at state-averaged orbitals, additional response terms are required; see
%\cref{sec:state_average}.
Only the part of $\Pi_{\mu \nu,\rho \sigma}$ symmetric with respect to interchange of the two orbital
pairs contributes to \cref{eq:e2_general}.
Therefore, we define
\begin{equation}
\overline{\Pi}_{\mu\nu,\rho\sigma}
= \frac{1}{2}
\left(
\Pi_{\mu\nu,\rho\sigma}
+ \Pi_{\rho\sigma,\mu\nu}
\right).
\label{eq:pair_symmetric_density}
\end{equation}
Because
$(\mu\nu|\rho\sigma)_\df = (\rho\sigma|\mu\nu)_\df$, the replacement
$\Pi \mapsto \overline{\Pi}$ leaves the energy unchanged.  In the remainder of
the derivation, $\Pi$ denotes this pair-symmetric density.

Substituting \cref{eq:df_eri} into \cref{eq:e2_general} gives
\begin{equation}
E
=  V_\mathrm{NN} +
\sum_{\mu \nu} h_{\mu\nu} \Pi_{\mu \nu}
+ \frac{1}{2}
\sum_{\mu \nu \rho \sigma} \sum_{PQ} \Pi_{\mu\nu,\rho\sigma}
J^P_{\mu\nu} N_{PQ} J^Q_{\rho\sigma}.
\label{eq:e2_df_raw}
\end{equation}

\section{First Derivative of the One-Electron Term}
We begin by computing the first derivative of the one-electron energy:
\begin{equation}
E_1 = \sum_{\mu \nu} h_{\mu\nu} \Pi_{\mu \nu}.
\end{equation}
We evaluate the derivative with respect to the parameter $x$ at constant one-electron density, obtaining:
\begin{equation}
E_1^{\pert} = \sum_{\mu \nu} h_{\mu\nu}^{\pert} \Pi_{\mu \nu}.
\end{equation}
This quantity is computed in Forte2 using the functions \texttt{kinetic\_deriv} and \texttt{nuclear\_deriv}, which take $\Pi_{\mu \nu}$ and return a vector of nuclear gradients of dimension $3M$.

\section{First Derivative of the Two-Electron Term}

Taking the first differential of \cref{eq:e2_df_raw} at fixed pair density gives three terms:
\begin{align}
E_2^{\pert}
&=
\frac{1}{2}
\Pi_{\mu\nu,\rho\sigma}
(J^P_{\mu\nu})^{\pert} N_{PQ} J^Q_{\rho\sigma}
\nonumber\\
&\quad+
\frac{1}{2}
\Pi_{\mu\nu,\rho\sigma}
J^P_{\mu\nu} (N_{PQ})^{\pert} J^Q_{\rho\sigma}
\nonumber\\
&\quad+
\frac{1}{2}
\Pi_{\mu\nu,\rho\sigma}
J^P_{\mu\nu} N_{PQ} (J^Q_{\rho\sigma})^{\pert}.
\label{eq:de2_three_terms}
\end{align}
The first and third terms are equal after exchanging the orbital pair labels
$(\mu\nu) \leftrightarrow (\rho\sigma)$ and using the symmetry of $\Pi$ and
$N$.  Thus
\begin{equation}
E_2^{\pert}
=
\Pi_{\mu\nu,\rho\sigma}
(J^P_{\mu\nu})^{\pert} N_{PQ} J^Q_{\rho\sigma}
+
\frac{1}{2}
\Pi_{\mu\nu,\rho\sigma}
J^P_{\mu\nu} (N_{PQ})^{\pert} J^Q_{\rho\sigma}
\label{eq:de2_after_pair_sym}
\end{equation}

Define the metric-applied three-center tensor
\begin{equation}
Z^P_{\mu\nu}
\equiv
N_{PQ} J^Q_{\mu\nu}.
\label{eq:z_ao}
\end{equation}
The three-center part of the differential is then
\begin{equation}
E_{2,J}^{\pert}
=
\Pi_{\mu\nu,\rho\sigma}
Z^P_{\rho\sigma}
(J^P_{\mu\nu})^{\pert}.
\label{eq:de2_j_preweight}
\end{equation}

The differential of the inverse metric follows from
\begin{equation}
N M = I.
\end{equation}
Differentiating this identity gives
\begin{equation}
N^{\pert} M + N M^{\pert} = 0,
\end{equation}
and multiplication by $N$ on the right yields
\begin{equation}
N^{\pert} = -NM^{\pert}N.
\label{eq:dinv_metric}
\end{equation}
In index notation,
\begin{equation}
N^{\pert}_{PQ}
=
-N_{PR}M^{\pert}_{RS}N_{SQ}.
\end{equation}
Substitution into the metric term in \cref{eq:de2_after_pair_sym} gives
\begin{align}
E_{2,M}^{\pert}
&=
-\frac{1}{2}
\Pi_{\mu\nu,\rho\sigma}
J^P_{\mu\nu}
N_{PR}
M_{RS}^{\pert}
N_{SQ}
J^Q_{\rho\sigma}
\nonumber\\
&=
-\frac{1}{2}
\Pi_{\mu\nu,\rho\sigma}
Z^R_{\mu\nu}
Z^S_{\rho\sigma}
M_{RS}^{\pert}.
\label{eq:de2_m_preweight}
\end{align}

Equations \cref{eq:de2_j_preweight,eq:de2_m_preweight} show that the full
explicit two-electron differential can be written as
\begin{equation}
E_2^{\pert}
=
W^P_{\mu\nu} (J^P_{\mu\nu})^{\pert}
+
W_{PQ} M_{PQ}^{\pert},
\label{eq:de2_weight_form}
\end{equation}
where the three-center derivative weight is
\begin{equation}
W^P_{\mu\nu}
=
\Pi_{\mu\nu,\rho\sigma}
Z^P_{\rho\sigma},
\label{eq:w3_ao}
\end{equation}
and the metric derivative weight is
\begin{equation}
W_{PQ}
=
-\frac{1}{2}
\Pi_{\mu\nu,\rho\sigma}
Z^P_{\mu\nu}
Z^Q_{\rho\sigma}.
\label{eq:w2_aux}
\end{equation}
Because $\Pi$ is pair symmetric, $W^M_{PQ}=W^M_{QP}$.  If an implementation
constructs an unsymmetrized pair density, \cref{eq:w3_ao,eq:w2_aux} should be
evaluated with $\overline{\Pi}$ from \cref{eq:pair_symmetric_density}.

For a nuclear coordinate $X$, the explicit two-electron gradient contribution is
\begin{equation}
E_2^{\pert}
=
W_{PQ} M^{\pert}_{PQ}
+
W^P_{\mu\nu} (J^{P}_{\mu\nu})^{\pert}
.
\label{eq:e2_gradient_final_ao}
\end{equation}
Written with integral notation, this is
\begin{equation}
E_2^{\pert}
=
W_{PQ} (P|Q)^{\pert}
+
W^P_{\mu\nu} (P|\mu\nu)^{\pert}
.
\label{eq:e2_gradient_integral_notation}
\end{equation}
This is the central result for implementation.  The first requires the
first derivatives of two-center two-electron Coulomb metric integrals.
The second contraction requires the first derivatives of three-center two-electron integrals.
In Forte2, we denote $W_{PQ}$ as $W_2$ and $W^P_{\mu\nu}$ as $W_3$.

%\section{Why Use the Raw Metric Form?}
%
%Many density-fitted algorithms store an orthonormalized three-index tensor
%\begin{equation}
%B^A_{\mu\nu} = X_{AP}J^P_{\mu\nu},
%\end{equation}
%where the metric factor $X$ satisfies
%\begin{equation}
%X^T X = N = M^{-1}.
%\end{equation}
%The DF integral may then be written
%\begin{equation}
%(\mu\nu|\rho\sigma)_\df
%= B^A_{\mu\nu}B^A_{\rho\sigma}.
%\end{equation}
%Although this form is useful for Fock builds, differentiating $B$ directly is
%not the preferred route:
%\begin{equation}
%(B^A_{\mu\nu})^{[X]}
%=
%X^{[X]}_{AP}J^P_{\mu\nu}
%+
%X_{AP}J^{P,[X]}_{\mu\nu}.
%\end{equation}
%The derivative $X^{[X]}$ depends on the particular square-root representation of
%the inverse metric and becomes especially delicate when the metric is
%orthogonalized or truncated.  The raw-metric derivation avoids differentiating
%$X$ entirely.  All metric dependence enters through the invariant identity
%\begin{equation}
%N^{[X]} = -N M^{[X]} N,
%\end{equation}
%which leads to the single metric derivative contraction in
%\cref{eq:e2_gradient_final_ao}.  This also makes the formula independent of
%whether $N$ is applied by Cholesky solves, an eigenvalue-truncated inverse, or
%another stable linear solve.

\section{Molecular-Orbital Form}

In multiconfigurational methods the pair density is naturally available in the
molecular-orbital basis.  Let $C_{\mu p}$ be the molecular-orbital coefficient
matrix and define
\begin{equation}
J^P_{pq}
=
\sum_{\mu\nu} C^*_{\mu p} J^P_{\mu\nu} C_{\nu q}.
\label{eq:j_mo}
\end{equation}
The metric-applied tensor in the MO basis is
\begin{equation}
Z^P_{pq}
=\sum_{Q}
N_{PQ}J^Q_{pq}.
\label{eq:z_mo}
\end{equation}
If the two-electron energy is
\begin{equation}
E_2
=
\frac{1}{2}
\sum_{pqrs}^\mathbb{G}\Pi_{pq,rs}
J^P_{pq}N_{PQ}J^Q_{rs},
\label{eq:e2_mo}
\end{equation}
where the pair density is related to the spin-free 2-RDM as:
\begin{equation}
\Pi_{pq,rs} = \Gamma_{pr,qs},
\end{equation}
then the MO-basis derivative weights are
\begin{align}
W^P_{pq}
&=
\sum_{rs}^\mathbb{G} \Pi_{pq,rs} Z^P_{rs} = \sum_{rs}^\mathbb{G}  \Gamma_{pr,qs} Z^P_{rs},
\label{eq:w3_mo}\\
W_{PQ}
&=
-\frac{1}{2}
\sum_{pqrs}^\mathbb{G}
\Pi_{pq,rs}Z^P_{pq}Z^Q_{rs}
=
-\frac{1}{2}
\sum_{pqrs}^\mathbb{G}
\Gamma_{pr,qs} Z^P_{pq} Z^Q_{rs}.
\label{eq:w2_mo}
\end{align}
The metric weight is already in the auxiliary basis and needs no orbital
transformation.  The three-center weight must be transformed back to the AO
basis for contraction with AO derivative integrals:
\begin{equation}
W^P_{\mu\nu}
=
C^*_{\mu p} W^P_{pq} C_{\nu q}.
\label{eq:w3_backtransform}
\end{equation}
For real orbitals this reduces to the usual transformation
$W^P_{\mu\nu}=C_{\mu p}W^P_{pq}C_{\nu q}$.

Using \cref{eq:w3_backtransform}, the MO-density form gives the same AO
gradient expression as \cref{eq:e2_gradient_final_ao}:
\begin{equation}
E_2^{\pert}
=
\sum_{\mu \nu}\left(\sum_{pq}^\mathbb{G}C^*_{\mu p} W^P_{pq} C_{\nu q}\right)(J^{P}_{\mu\nu})^{\pert}
+
\sum_{PQ}
W_{PQ}M^{\pert}_{PQ}.
\label{eq:e2_gradient_mo_to_ao}
\end{equation}

\section{SCF Specializations and Required Quantities}
\label{sec:scf_requirements}

For any single-determinant SCF reference, the two-electron density is generated
from the one-particle density by \cref{eq:determinant_spinorbital_2rdm}.  The
DF gradient formulas therefore require only the occupied orbitals, the
spin-block one-particle densities, and the metric-applied three-center tensor
$Z$.

A complete analytic nuclear gradient additionally requires the one-electron
and overlap/Lagrangian contributions.  In the notation of the previous
sections, the method-dependent quantities needed for a full SCF gradient are:
\begin{enumerate}
\item the one-particle density matrix for $h^{\pert}$;
\item the pair density $\Pi_{\mu\nu,\rho\sigma}$ or an equivalent factorization for the
      two-electron derivative terms in \cref{eq:e2_gradient_final_ao};
\item the orbital Lagrange multiplier, often represented as an
      energy-weighted density, for the overlap derivative contribution;
\item the nuclear repulsion derivative.
\end{enumerate}
Only item 2 is specific to the density-fitted two-electron derivation in this
note, but all four are needed to produce a full gradient.

\subsection{RHF}

For RHF, let $C_{\mu i}$ be the occupied spatial orbitals and define the
density matrix
\begin{equation}
D_{\mu\nu}
=
\sum_i^{\mathbb{O}}
C^*_{\mu i}C_{\nu i}.
\end{equation}
The spin-free charge density is
\begin{equation}
P_{\mu\nu}=2D_{\mu\nu}.
\end{equation}
For a RHF state, we have that:
\begin{equation}
\Pi_{\mu\nu,\rho\sigma}^{\mathrm{RHF}}
=
P_{\mu\nu}P_{\rho\sigma}
-
\frac{1}{2}P_{\mu\sigma}P_{\rho\nu}.
\end{equation}

This gives:
\begin{equation}
W^P_{\mu\nu}(\mathrm{RHF})
=
P_{\mu\nu} (\sum_{\rho\sigma}P_{\rho\sigma}
Z^P_{\rho\sigma})
-\frac{1}{2}\sum_{\sigma}P_{\mu\sigma}(\sum_{\rho} P_{\rho\nu}Z^P_{\rho\sigma}),
\end{equation}
which we express as:
\begin{equation}
W^P_{\mu\nu}(\mathrm{RHF})
=
P_{\mu\nu} R^P
-2\sum_{ij} C_{\mu i}^* C_{\nu j} Z^P_{ji},
\end{equation}
where
\begin{equation}
Z^P_{ij} = \sum_{\mu\nu} C_{\mu i}^* C_{\nu j} Z^P_{\mu\nu},
\end{equation}
and
\begin{equation}
R^P = \sum_{\rho\sigma}P_{\rho\sigma}
Z^P_{\rho\sigma}.
\end{equation}

The metric derivative weight is
\begin{equation}
W_{PQ}
=
-\frac{1}{2}
(\sum_{\mu \nu} P_{\mu\nu}Z^P_{\mu\nu})(\sum_{\rho\sigma} P_{\rho\sigma}
Z^Q_{\rho\sigma})
+
\frac{1}{4}
\sum_{\nu\sigma}(\sum_\mu P_{\mu\sigma}Z^P_{\mu\nu}) (\sum_\rho P_{\rho\nu}
Z^Q_{\rho\sigma}).
\label{eq:rhf_w2_aux}
\end{equation}
which we express as:
\begin{equation}
W_{PQ}
=
-\frac{1}{2}
R^P R^Q
+
\sum_{ij}^\mathbb{O} Z^P_{ij} Z^Q_{ji}.
\label{eq:rhf_w2_aux}
\end{equation}

\subsection{UHF}

For UHF, the alpha and beta occupied orbitals define two collinear spin
densities
\begin{align}
D^\alpha_{\mu\nu}
&=
\sum_i^{\mathbb{O}_\alpha}
C^{\alpha *}_{\mu i}C^\alpha_{\nu i},
\\
D^\beta_{\mu\nu}
&=
\sum_i^{\mathbb{O}_\beta}
C^{\beta *}_{\mu i}C^\beta_{\nu i}.
\end{align}
The charge density is
\begin{equation}
P_{\mu\nu}=D^\alpha_{\mu\nu}+D^\beta_{\mu\nu},
\end{equation}
and the pair density is
\begin{equation}
\Pi_{\mu\nu,\rho\sigma}^{\mathrm{UHF}}
=
P_{\mu\nu}P_{\rho\sigma}
-
D^\alpha_{\mu\sigma}D^\alpha_{\rho\nu}
-
D^\beta_{\mu\sigma}D^\beta_{\rho\nu}.
\label{eq:uhf_pair_density}
\end{equation}

Define
\begin{equation}
Z^{P,\alpha}_{ij} = \sum_{\mu\nu} C_{\mu i}^{\alpha*} C_{\nu j}^{\alpha} Z^P_{\mu\nu},
\end{equation}
and
\begin{equation}
Z^{P,\beta}_{ij} = \sum_{\mu\nu} C_{\mu i}^{\beta*} C_{\nu j}^{\beta} Z^P_{\mu\nu},
\end{equation}
as well as:
\begin{equation}
R^P = \sum_{\rho\sigma}P_{\rho\sigma}
Z^P_{\rho\sigma}.
\end{equation}

This gives:
\begin{equation}
W^P_{\mu\nu}(\mathrm{UHF})
=
P_{\mu\nu} R^P
-\sum_{\sigma\rho}D^{\alpha}_{\mu\sigma} D^{\alpha}_{\rho\nu}Z^P_{\rho\sigma}
-\sum_{\sigma\rho}D^{\beta}_{\mu\sigma} D^{\beta}_{\rho\nu}Z^P_{\rho\sigma},
\end{equation}
which we express as:
\begin{equation}
W^P_{\mu\nu}(\mathrm{UHF})
=
P_{\mu\nu} R^P
-\sum_{ij}^{\mathbb{O}_\alpha} C^{\alpha*}_{\mu i} C^{\alpha}_{\nu j} Z^{P,\alpha}_{ji}
-\sum_{ij}^{\mathbb{O}_\beta} C^{\beta*}_{\mu i} C^{\beta}_{\nu j} Z^{P,\beta}_{ji}.
\label{eq:uhf_w3_aux}
\end{equation}

The metric derivative weight is
\begin{equation}
W_{PQ}
=
-\frac{1}{2}
R^P R^Q
+
\frac{1}{2}
\sum_{\mu\nu\rho\sigma}D^{\alpha}_{\mu\sigma}D^{\alpha}_{\rho\nu}Z^P_{\mu\nu}Z^Q_{\rho\sigma}
+
\frac{1}{2}
\sum_{\mu\nu\rho\sigma}D^{\beta}_{\mu\sigma}D^{\beta}_{\rho\nu}Z^P_{\mu\nu}Z^Q_{\rho\sigma}.
\end{equation}
which we express as:
\begin{equation}
W_{PQ}
=
-\frac{1}{2}
R^P R^Q
+\frac{1}{2}
\sum_{ij}^\mathbb{O_\alpha} Z^{P,\alpha}_{ij} Z^{Q,\alpha}_{ji}
+\frac{1}{2}
\sum_{ij}^\mathbb{O_\beta} Z^{P,\beta}_{ij} Z^{Q,\beta}_{ji}.
\label{eq:uhf_w2_aux}
\end{equation}


%
%\subsection{ROHF}
%
%ROHF uses a single spatial orbital set but different alpha and beta occupations.
%For the explicit two-electron gradient, ROHF is algebraically the same
%collinear determinant as UHF:
%\begin{equation}
%D^\alpha_{\mu\nu}
%=
%\sum_{i\in \alpha\ \mathrm{occ}} C^*_{\mu i}C_{\nu i},
%\qquad
%D^\beta_{\mu\nu}
%=
%\sum_{i\in \beta\ \mathrm{occ}} C^*_{\mu i}C_{\nu i}.
%\end{equation}
%The same pair density in \cref{eq:uhf_pair_density} therefore generates the
%two-electron derivative weights.  In terms of the usual closed-shell and
%open-shell densities,
%\begin{align}
%D^{\mathrm{closed}}_{\mu\nu}
%&=
%\sum_{i\in \mathrm{doubly\ occ}} C^*_{\mu i}C_{\nu i},
%\\
%D^{\mathrm{open}}_{\mu\nu}
%&=
%\sum_{i\in \mathrm{singly\ occ}} C^*_{\mu i}C_{\nu i},
%\end{align}
%one commonly has, for the high-spin component,
%\begin{equation}
%D^\alpha=D^{\mathrm{closed}}+D^{\mathrm{open}},
%\qquad
%D^\beta=D^{\mathrm{closed}}.
%\end{equation}
%The quantities needed for the two-electron DF part are therefore the common
%orbital coefficients, the alpha and beta occupation projectors, the resulting
%$D^\alpha$, $D^\beta$, $P$, and the DF weights.
%
%The overlap derivative contribution is the subtle part of ROHF.  Because the
%ROHF canonical Fock matrix is not unique, the energy-weighted density must be
%constructed from the ROHF Lagrangian or from the same canonicalization
%convention used by the SCF procedure.  It is not sufficient to reuse the RHF
%or UHF occupied orbital energy formula without checking consistency with the
%chosen ROHF effective Fock matrix.
%
%\subsection{GHF and Complex GHF}
%
%For GHF, each occupied spinor has alpha and beta spatial components,
%\begin{equation}
%\psi_i(\mathbf{x})
%=
%\sum_\mu
%C^\alpha_{\mu i}\chi_\mu(\mathbf{r})\alpha
%+
%\sum_\mu
%C^\beta_{\mu i}\chi_\mu(\mathbf{r})\beta .
%\end{equation}
%The spin-block one-particle densities are
%\begin{equation}
%\gamma^{\omega\tau}_{\mu\nu}
%=
%\sum_i^{\mathrm{occ}}
%C^{\omega *}_{\mu i}C^\tau_{\nu i},
%\qquad
%\omega,\tau\in\{\alpha,\beta\}.
%\label{eq:ghf_spin_block_density}
%\end{equation}
%These density blocks may be complex.  Hermiticity of the full spinor density
%implies
%\begin{equation}
%\gamma^{\tau\omega}_{\nu\mu}
%=
%\left(\gamma^{\omega\tau}_{\mu\nu}\right)^*.
%\end{equation}
%The charge density entering the Coulomb part is still only the spin trace,
%\begin{equation}
%P_{\mu\nu}
%=
%\gamma^{\alpha\alpha}_{\mu\nu}
%+
%\gamma^{\beta\beta}_{\mu\nu}.
%\end{equation}
%The exchange part, however, depends on all four spin blocks:
%\begin{equation}
%\Pi_{\mu\nu,\rho\sigma}^{\mathrm{GHF}}
%=
%P_{\mu\nu}P_{\rho\sigma}
%-
%\sum_{\omega\tau}
%\gamma^{\omega\tau}_{\mu\sigma}
%\gamma^{\tau\omega}_{\rho\nu}.
%\label{eq:ghf_pair_density}
%\end{equation}
%The off-diagonal spin densities
%$\gamma^{\alpha\beta}$ and $\gamma^{\beta\alpha}$ are what distinguish GHF
%from a collinear UHF determinant.
%
%For the two-electron DF gradient, complex GHF therefore needs:
%\begin{enumerate}
%\item occupied complex spinor coefficients $C^\alpha_{\mu i}$ and
%      $C^\beta_{\mu i}$;
%\item all four spin-block densities
%      $\gamma^{\alpha\alpha}$, $\gamma^{\alpha\beta}$,
%      $\gamma^{\beta\alpha}$, and $\gamma^{\beta\beta}$;
%\item the spin-traced charge density $P$ for the Coulomb part;
%\item all spin-block exchange contractions from
%      \cref{eq:ghf_pair_density};
%\item complex weights $W^P_{\mu\nu}$ and $W^M_{PQ}$, followed by taking the
%      real part of the final scalar gradient component.
%\end{enumerate}
%The derivative integrals $(P|\mu\nu)^{[X]}$ and $(P|Q)^{[X]}$ are real for a
%real scalar AO and auxiliary basis, but the density weights can be complex.
%The exact gradient is real; any residual imaginary part in the contraction is
%numerical noise if the density and Lagrangian are Hermitian and stationary.
%
%The one-electron and overlap terms for GHF must be evaluated in the spinor AO
%basis.  For a nonrelativistic scalar Hamiltonian, one-electron and overlap
%integral derivatives are spin diagonal, but the orbital Lagrangian and density
%are spinor quantities.  In a complex GHF implementation, the energy-weighted
%density must be built with complex conjugation consistent with the occupied
%spinor coefficients and the converged spinor Fock matrix.
%
%\section{CASSCF Pair Density}
%
%For CASSCF and GASSCF, the relevant two-electron density is the effective
%spin-free pair density over the internal orbitals that enter the CASSCF
%energy expression.  Introduce the spin-free excitation operators
%\begin{equation}
%\hat{E}_{pq}
%=
%\sum_{\sigma}
%\hat{a}^\dagger_{p\sigma}\hat{a}_{q\sigma}.
%\end{equation}
%The spin-free two-electron Hamiltonian can be written as
%\begin{equation}
%\hat{V}
%=
%\frac{1}{2}
%(pq|rs)
%\left(
%\hat{E}_{pq}\hat{E}_{rs}
%-\delta_{qr}\hat{E}_{ps}
%\right).
%\label{eq:spinfree_v}
%\end{equation}
%Therefore the CASSCF effective pair density is
%\begin{equation}
%\Pi_{pq,rs}
%=
%\left\langle
%\hat{E}_{pq}\hat{E}_{rs}
%-\delta_{qr}\hat{E}_{ps}
%\right\rangle.
%\label{eq:casscf_pair_density_operator}
%\end{equation}
%With this definition,
%\begin{equation}
%E_2^\mathrm{CASSCF}
%=
%\frac{1}{2}
%\Pi_{pq,rs}
%(pq|rs)_\df.
%\label{eq:casscf_e2_pair_density}
%\end{equation}
%
%The indices $p,q,r,s$ in \cref{eq:casscf_pair_density_operator} run over the
%orbital space included in the CASSCF two-electron energy expression.  In the
%usual closed-core CASSCF partition this includes core and active orbitals.  The
%closed-shell core part of \cref{eq:casscf_pair_density_operator} reproduces the
%standard core-core Coulomb and exchange energy, while mixed core-active and
%active-active blocks are determined by the active-space one- and two-particle
%RDMs.
%
%Equivalently, an implementation may avoid explicitly materializing the full
%core-plus-active pair density by accumulating the same weights
%\cref{eq:w3_mo,eq:w2_mo} from the separate CASSCF energy pieces:
%\begin{equation}
%E_2^\mathrm{CASSCF}
%=
%E_2^{\core\core}
%+
%E_2^{\core\act}
%+
%E_2^{\act\act}.
%\end{equation}
%The total weights are additive:
%\begin{align}
%W^P_{pq}
%&=
%W^{P,\core\core}_{pq}
%+
%W^{P,\core\act}_{pq}
%+
%W^{P,\act\act}_{pq},
%\\
%W^M_{PQ}
%&=
%W^{M,\core\core}_{PQ}
%+
%W^{M,\core\act}_{PQ}
%+
%W^{M,\act\act}_{PQ}.
%\end{align}
%This blockwise implementation is algebraically identical to forming the full
%$\Pi_{pq,rs}$ in \cref{eq:casscf_pair_density_operator}, as long as all blocks
%use the same integral index convention.
%
%\subsection{Active-Space Contribution}
%
%For the purely active part, let $u,v,w,x$ denote active orbitals and let
%$\Pi_{uv,wx}^{\act}$ be the active-space effective pair density in the same
%index convention as \cref{eq:casscf_pair_density_operator}.  Then
%\begin{equation}
%E_2^{\act\act}
%=
%\frac{1}{2}
%\Pi_{uv,wx}^{\act}
%J^P_{uv}N_{PQ}J^Q_{wx}.
%\end{equation}
%The corresponding weights are
%\begin{align}
%W^{P,\act\act}_{uv}
%&=
%\Pi_{uv,wx}^{\act}Z^P_{wx},
%\\
%W^{M,\act\act}_{PQ}
%&=
%-\frac{1}{2}
%\Pi_{uv,wx}^{\act}Z^P_{uv}Z^Q_{wx}.
%\end{align}
%These are the direct analogues of \cref{eq:w3_mo,eq:w2_mo}.
%
%\subsection{Core and Core-Active Contributions}
%
%The closed-shell core contribution can also be generated from
%\cref{eq:casscf_pair_density_operator}.  For example, for core orbitals $i,j$
%only, the effective pair density is
%\begin{equation}
%\Pi_{ij,kl}^{\core}
%=
%4\delta_{ij}\delta_{kl}
%-
%2\delta_{il}\delta_{jk}.
%\label{eq:core_pair_density}
%\end{equation}
%Substitution into
%$E_2=(1/2)\Pi_{ij,kl}(ij|kl)$ gives
%\begin{equation}
%E_2^{\core\core}
%=
%2(ii|jj) - (ij|ji),
%\end{equation}
%with summation over repeated core indices.  Mixed core-active blocks are
%obtained from the same operator definition and depend on the active one-particle
%RDM.  In practice, the least ambiguous implementation is to either:
%\begin{enumerate}
%\item build the complete effective internal pair density
%      $\Pi_{pq,rs}$ from \cref{eq:casscf_pair_density_operator}, or
%\item build the three-center and metric weights by differentiating the same
%      block energy expressions used by the CASSCF energy code.
%\end{enumerate}
%Both routes must yield the same final weights in
%\cref{eq:w3_mo,eq:w2_mo}.
%
%\section{State-Averaged CASSCF}
%\label{sec:state_average}
%
%For state-averaged CASSCF, the optimized objective is
%\begin{equation}
%E_\sa
%=
%\sum_I w_I E_I,
%\qquad
%\sum_I w_I = 1,
%\qquad
%w_I \geq 0,
%\end{equation}
%where $I$ labels the CI roots included in the state average.  The two-electron
%part is
%\begin{equation}
%E_{2,\sa}
%=
%\sum_I w_I E_{2,I}
%=
%\frac{1}{2}
%\left(
%\sum_I w_I \Pi^I_{pq,rs}
%\right)
%(pq|rs)_\df.
%\end{equation}
%Therefore define the state-averaged pair density
%\begin{equation}
%\Pi^\sa_{pq,rs}
%=
%\sum_I w_I \Pi^I_{pq,rs}.
%\label{eq:sa_pair_density}
%\end{equation}
%The state-averaged two-electron gradient contribution is obtained by using
%$\Pi^\sa$ in the same equations derived above:
%\begin{align}
%W^{P,\sa}_{pq}
%&=
%\Pi^\sa_{pq,rs}Z^P_{rs},
%\\
%W^{M,\sa}_{PQ}
%&=
%-\frac{1}{2}
%\Pi^\sa_{pq,rs}Z^P_{pq}Z^Q_{rs},
%\\
%E_{2,\sa}^{[X]}
%&=
%\left(C^*_{\mu p}W^{P,\sa}_{pq}C_{\nu q}\right)
%J^{P,[X]}_{\mu\nu}
%+
%W^{M,\sa}_{PQ}M^{[X]}_{PQ}.
%\end{align}
%
%This expression is valid for the gradient of the state-averaged objective at a
%state-averaged stationary point.  If instead one wants the gradient of a single
%root $K$ evaluated at orbitals optimized for a different state-averaged
%objective, the target root energy is not stationary with respect to those
%orbitals.  The explicit integral derivative contribution is still obtained by
%using $\Pi^K$, but the full analytic gradient also requires orbital and CI
%response, commonly expressed through a CASSCF response or Z-vector equation.
%
%\section{CASSCF Z-Vector Equations}
%\label{sec:casscf_zvector}
%
%The preceding sections give the explicit derivative of the two-electron
%integrals at fixed molecular orbitals and fixed CI coefficients.  For a fully
%variational, state-specific CASSCF energy, the fixed-wave-function expression
%is already sufficient apart from the usual overlap/Lagrangian term: the energy
%is stationary with respect to all nonredundant orbital rotations and CI
%parameters.  However, a Z-vector is needed whenever the target energy is not
%the same functional that defined the optimized orbitals and CI coefficients.
%The most common examples are:
%\begin{enumerate}
%\item the gradient of an individual root at state-averaged CASSCF orbitals;
%\item gradients with orbital spaces held fixed or frozen, such as frozen-core
%      orbitals excluded from the CASSCF orbital optimization;
%\item any property or post-CASSCF Lagrangian whose density is not stationary
%      with respect to the CASSCF variables.
%\end{enumerate}
%
%\subsection{Variables and Stationarity Conditions}
%
%Let $y$ collect all independent CASSCF wave-function variables:
%\begin{equation}
%y =
%\begin{pmatrix}
%\boldsymbol{\kappa} \\
%\boldsymbol{\theta}
%\end{pmatrix},
%\end{equation}
%where $\boldsymbol{\kappa}$ are the nonredundant anti-Hermitian orbital
%rotations and $\boldsymbol{\theta}$ parameterizes the CI coefficients in the
%active space.  In a real spin-free implementation, the nonredundant orbital
%rotations are the unique lower-triangular elements connecting different orbital
%spaces, for example core-active, core-virtual, active-virtual, and, for GASSCF,
%inter-GAS rotations.  In a complex implementation, the Lagrangian is real and
%the transpose operations below should be understood as Hermitian adjoints in
%the real-valued parameterization of complex variables.
%
%Let $E_R(y)$ be the reference objective used to determine the CASSCF wave
%function.  For state-specific CASSCF, $E_R$ is the energy of that state.  For
%state-averaged CASSCF,
%\begin{equation}
%E_R = E_\sa = \sum_I w_I E_I.
%\end{equation}
%At convergence, the reference stationarity equations are
%\begin{equation}
%g_m^R(y) \equiv \frac{\partial E_R}{\partial y_m} = 0,
%\qquad
%m \in \mathcal{V},
%\label{eq:casscf_stationarity}
%\end{equation}
%where $\mathcal{V}$ is the set of optimized CASSCF variables.  If some orbital
%rotations are excluded, for example rotations involving frozen core orbitals,
%they are not included in $\mathcal{V}$ unless additional constraints are added
%as described in \cref{subsec:frozen_core_zvector}.
%
%Now let $E_T(y)$ be the target energy whose gradient is desired.  The total
%derivative with respect to a nuclear coordinate $X$ is
%\begin{equation}
%\frac{\dd E_T}{\dd X}
%=
%E_T^{[X]}
%+
%\frac{\partial E_T}{\partial y_m} y_m^{[X]}.
%\label{eq:target_total_derivative}
%\end{equation}
%Differentiating the stationarity conditions gives the response equation
%\begin{equation}
%A^R_{mn} y_n^{[X]} + g_m^{R,[X]} = 0,
%\qquad
%A^R_{mn}
%=
%\frac{\partial g_m^R}{\partial y_n}
%=
%\frac{\partial^2 E_R}{\partial y_m \partial y_n}.
%\label{eq:casscf_response_equation}
%\end{equation}
%Solving \cref{eq:casscf_response_equation} separately for every nuclear
%coordinate is avoidable.  Define the Z-vector by
%\begin{equation}
%\left(A^R\right)^T \mathbf{z}
%=
%-\mathbf{b},
%\qquad
%b_m =
%\frac{\partial E_T}{\partial y_m}.
%\label{eq:zvector_equation}
%\end{equation}
%Then \cref{eq:target_total_derivative} becomes
%\begin{equation}
%\frac{\dd E_T}{\dd X}
%=
%E_T^{[X]}
%+
%z_m g_m^{R,[X]}.
%\label{eq:zvector_gradient}
%\end{equation}
%This is the basic CASSCF Z-vector result.  If the target is the same stationary
%functional as the reference, then $\mathbf{b}=0$ and $\mathbf{z}=0$.
%
%\subsection{Coupled Orbital-CI Form}
%
%Partitioning the variables into orbital rotations and CI parameters gives
%\begin{equation}
%\begin{pmatrix}
%A_{\kappa\kappa}^R & A_{\kappa c}^R \\
%A_{c\kappa}^R & A_{cc}^R
%\end{pmatrix}^T
%\begin{pmatrix}
%z^\kappa \\
%z^c
%\end{pmatrix}
%=
%-
%\begin{pmatrix}
%b^\kappa \\
%b^c
%\end{pmatrix}.
%\label{eq:coupled_casscf_zvector}
%\end{equation}
%The orbital right-hand side is the target orbital gradient,
%\begin{equation}
%b^\kappa_{pq}
%=
%\frac{\partial E_T}{\partial \kappa_{pq}},
%\end{equation}
%evaluated at the reference CASSCF orbitals and CI vectors.  The CI right-hand
%side is the target CI gradient.  For an individual state energy, the CI vector
%for that state is an eigenvector of the active-space Hamiltonian at fixed
%orbitals, so $b^c=0$.  Nevertheless, the CI part of the Z-vector can still be
%nonzero because the orbital and CI variables are coupled by
%$A_{\kappa c}^R$ and $A_{c\kappa}^R$.
%
%For a single root $I$ with CI vector $\mathbf{c}_I$, the CI stationarity
%condition may be written as the projected residual
%\begin{equation}
%\mathbf{r}_I
%=
%Q_I
%\left(
%\mathbf{H}^{\act} - E_I \mathbf{1}
%\right)
%\mathbf{c}_I
%=
%0,
%\label{eq:ci_projected_residual}
%\end{equation}
%where
%\begin{equation}
%Q_I = \mathbf{1} - \mathbf{c}_I \mathbf{c}_I^\dagger
%\end{equation}
%projects out the normalization direction.  The CI-CI block of the CASSCF
%Hessian is therefore the projected active-space Hamiltonian shift
%\begin{equation}
%A^R_{cc,I}
%=
%Q_I
%\left(
%\mathbf{H}^{\act} - E_I \mathbf{1}
%\right)
%Q_I
%\label{eq:ci_ci_hessian}
%\end{equation}
%for a state-specific root.  In state-averaged calculations, the CI residuals
%are written for all roots in the average, with root orthogonality constraints
%included in the projector.  Near-degenerate roots should be treated in the
%state subspace rather than by assigning arbitrary individual-root phases.
%
%The orbital-CI coupling blocks are derivatives of the CI residuals with respect
%to orbital rotations and derivatives of the orbital gradient with respect to
%CI parameters:
%\begin{align}
%A^R_{c\kappa}
%&=
%\frac{\partial \mathbf{r}}{\partial \boldsymbol{\kappa}},
%\\
%A^R_{\kappa c}
%&=
%\frac{\partial \mathbf{g}^{R,\kappa}}{\partial \mathbf{c}}.
%\end{align}
%In practice, these blocks are most often applied as sigma-vector operations
%rather than formed explicitly.
%
%\subsection{Evaluation With Relaxed Density Matrices}
%
%After solving \cref{eq:coupled_casscf_zvector}, the gradient may be evaluated
%directly from \cref{eq:zvector_gradient}.  Equivalently, one can absorb the
%Z-vector terms into relaxed one- and two-particle density matrices and then use
%the same derivative-integral formulas as before.
%
%For the CI part, introduce a response CI vector $\mathbf{z}^c_I$ for root $I$,
%orthogonal to $\mathbf{c}_I$.  The CI-response corrections to the active-space
%spin-free RDMs have the form
%\begin{align}
%\gamma^{Z,c}_{uv}
%&=
%2\,\mathrm{Re}
%\left[
%\langle \mathbf{z}^c_I |
%\hat{E}_{uv}
%| \mathbf{c}_I \rangle
%\right],
%\label{eq:ci_z_1rdm}\\
%\Gamma^{Z,c}_{uv,wx}
%&=
%2\,\mathrm{Re}
%\left[
%\left\langle \mathbf{z}^c_I \left|
%\hat{E}_{uv}\hat{E}_{wx}
%-
%\delta_{vw}\hat{E}_{ux}
%\right| \mathbf{c}_I \right\rangle
%\right],
%\label{eq:ci_z_2rdm}
%\end{align}
%with the obvious weighted sum over roots for a state-averaged Lagrangian.  For
%complex wave functions, the real part is taken after the bra-ket contraction.
%
%The orbital Z-vector contributes through the derivative of the orbital
%stationarity equations, $z^\kappa_{pq}g_{pq}^{R,\kappa,[X]}$.  This term can be
%implemented either by applying derivative integrals to the corresponding
%orbital-response density contributions, or by evaluating the contracted
%Lagrangian derivative directly.  In either approach, the final two-electron
%gradient can still be written in the density-fitted weight form
%\begin{equation}
%E_2^{[X]}
%=
%\overline{W}^P_{\mu\nu}(P|\mu\nu)^{[X]}
%+
%\overline{W}^M_{PQ}(P|Q)^{[X]},
%\end{equation}
%but the weights are built from relaxed densities:
%\begin{align}
%\overline{\gamma}
%&=
%\gamma^T + \gamma^Z,
%\\
%\overline{\Pi}
%&=
%\Pi^T + \Pi^Z.
%\end{align}
%Here the superscript $T$ denotes the target-state or target-objective density,
%and the superscript $Z$ denotes all CI and orbital-response contributions from
%the Z-vector Lagrangian.
%
%\subsection{Frozen Core Orbitals}
%\label{subsec:frozen_core_zvector}
%
%Frozen core orbitals require extra bookkeeping because they are included as
%doubly occupied orbitals in the energy but excluded from the optimized CASSCF
%orbital-rotation space.  Let the contiguous orbital spaces be
%\begin{equation}
%\{p\}
%=
%\{f\}
%\oplus
%\{i\}
%\oplus
%\{u\}
%\oplus
%\{a\}
%\oplus
%\{r\},
%\end{equation}
%where $f$ are frozen core orbitals, $i$ are optimized inactive core orbitals,
%$u$ are active orbitals, $a$ are optimized virtual orbitals, and $r$ are frozen
%virtual orbitals if present.  In the Forte2 convention, the inactive doubly
%occupied space used in the CASSCF energy contains both $f$ and $i$, while
%nonredundant CASSCF orbital rotations are normally taken only among the
%non-frozen blocks, for example
%\begin{equation}
%\mathcal{R}_{\mathrm{CASSCF}}
%=
%\{u i,\; a i,\; a u\}
%\end{equation}
%for ordinary CASSCF, plus inter-GAS rotations for GASSCF.  Rotations involving
%$f$ are omitted from this set.
%
%The frozen core orbitals still contribute to the unrelaxed density.  In a
%spin-free closed-shell treatment,
%\begin{equation}
%\gamma_{ff'} \leftarrow 2\delta_{ff'},
%\end{equation}
%and the corresponding inactive-inactive and inactive-active pair-density blocks
%must be included exactly as for any other doubly occupied core orbital.  Thus
%the explicit DF derivative weights include the frozen core contribution even
%before any response treatment is considered.
%
%There are two distinct gradient conventions:
%\begin{enumerate}
%\item \emph{Constrained frozen-orbital gradient.}  The frozen core orbitals are
%      treated as fixed external orbitals defining a constrained variational
%      problem.  The Z-vector equation is solved only in the optimized CASSCF
%      variable space $\mathcal{V}=\mathcal{R}_{\mathrm{CASSCF}}\oplus
%      \mathcal{V}_{\mathrm{CI}}$.  Omitted rotations involving $f$ are not
%      stationarity constraints and are not relaxed.  This gives the derivative
%      of the constrained frozen-orbital energy.
%\item \emph{Parent-relaxed frozen-core gradient.}  The frozen core orbitals are
%      assumed to follow the parent mean-field orbitals as the geometry changes.
%      Then their response must be determined by the parent SCF stationarity
%      equations, and extra Z-vector multipliers are required for the parent
%      orbital-gradient constraints that define the frozen core subspace.
%\end{enumerate}
%
%For the parent-relaxed convention, define an additional set of orbital
%variables $\mathcal{F}$ that contains the nonredundant rotations involving the
%frozen core and the non-frozen orbital spaces whose response should be inherited
%from the parent reference.  Typical examples are $u f$ and $a f$ rotations.
%Rotations within a fully occupied inactive subspace, such as $i f$, are
%redundant for the CASSCF energy unless the frozen/correlated core distinction
%is used by a downstream method.
%
%The Lagrangian then contains both CASSCF and parent-SCF stationarity
%constraints:
%\begin{equation}
%\mathcal{L}
%=
%E_T
%+
%z_m^{\mathrm{CAS}} g_m^{\mathrm{CAS}}
%+
%z_\alpha^{\mathrm{fc}} g_\alpha^{\mathrm{parent}},
%\label{eq:frozen_core_lagrangian}
%\end{equation}
%where $m\in\mathcal{R}_{\mathrm{CASSCF}}\oplus\mathcal{V}_{\mathrm{CI}}$ and
%$\alpha\in\mathcal{F}$.  The generalized Z-vector equation is
%\begin{equation}
%\begin{pmatrix}
%\dfrac{\partial g_m^{\mathrm{CAS}}}{\partial y_n} &
%\dfrac{\partial g_m^{\mathrm{CAS}}}{\partial y_\beta} \\
%\dfrac{\partial g_\alpha^{\mathrm{parent}}}{\partial y_n} &
%\dfrac{\partial g_\alpha^{\mathrm{parent}}}{\partial y_\beta}
%\end{pmatrix}^T
%\begin{pmatrix}
%z_n^{\mathrm{CAS}} \\
%z_\beta^{\mathrm{fc}}
%\end{pmatrix}
%=
%-
%\begin{pmatrix}
%\dfrac{\partial E_T}{\partial y_m} \\
%\dfrac{\partial E_T}{\partial y_\alpha}
%\end{pmatrix}.
%\label{eq:frozen_core_zvector}
%\end{equation}
%Here $y_n$ are the optimized CASSCF variables and $y_\alpha$ are the
%parent-defined frozen-core rotation variables.  If $E_T$ is the state-averaged
%CASSCF objective, the first block of the right-hand side vanishes by CASSCF
%stationarity, but the frozen-core block
%\begin{equation}
%\frac{\partial E_T}{\partial y_\alpha}
%\end{equation}
%is generally nonzero because rotations involving frozen core orbitals were not
%optimized in the CASSCF calculation.
%
%After solving \cref{eq:frozen_core_zvector}, the gradient is evaluated as
%\begin{equation}
%\frac{\dd E_T}{\dd X}
%=
%E_T^{[X]}
%+
%z_m^{\mathrm{CAS}} g_m^{\mathrm{CAS},[X]}
%+
%z_\alpha^{\mathrm{fc}} g_\alpha^{\mathrm{parent},[X]}.
%\label{eq:frozen_core_gradient}
%\end{equation}
%The last term is the parent-relaxation correction for the frozen core orbitals.
%Operationally, it adds parent-SCF Z-density contributions to the one-electron,
%two-electron, and overlap derivative contractions.  The two-electron DF part
%can still be evaluated by the same weight equations, but using the density
%augmented by both the CASSCF Z-vector and frozen-core parent-response
%contributions.
%
%\subsection{Practical Solution Strategy}
%
%A practical CASSCF Z-vector implementation can proceed as follows:
%\begin{enumerate}
%\item Define the variable space: optimized orbital rotations, CI response
%      parameters, and, if using parent-relaxed frozen core orbitals, the
%      parent-defined frozen-core rotation variables.
%\item Build the right-hand side $\mathbf{b}$ from the target orbital and CI
%      gradients.  For the state-averaged objective itself, this is zero in the
%      optimized CASSCF variables.  For an individual state at state-averaged
%      orbitals, the orbital part is the individual-state orbital gradient.
%\item Implement Hessian-vector products for the CASSCF response matrix.  These
%      may be analytic sigma-vector operations or finite differences of the
%      existing CASSCF orbital and CI residual routines.  For frozen-core parent
%      relaxation, the parent block is the usual coupled-perturbed SCF orbital
%      Hessian.
%\item Project out redundant rotations, CI normalization directions, and root
%      rotations inside exactly degenerate averaged subspaces.  The linear
%      system is otherwise singular.
%\item Solve the adjoint equation \cref{eq:zvector_equation} or
%      \cref{eq:frozen_core_zvector} with an iterative linear solver.  The
%      orbital Hessian diagonal used in the CASSCF optimizer is a natural
%      preconditioner for the orbital block; CI energy gaps precondition the CI
%      block; parent orbital-energy differences precondition the parent-SCF
%      frozen-core block.
%\item Convert the converged Z-vector into relaxed one- and two-particle density
%      matrices, or evaluate the contracted Lagrangian derivative directly.
%\item Generate the density-fitted two-electron derivative weights from the
%      relaxed pair density and contract them with
%      $(P|\mu\nu)^{[X]}$ and $(P|Q)^{[X]}$.
%\end{enumerate}
%
%\section{DSRG-MRPT2 Density-Fitted Derivatives}
%\label{sec:dsrg_mrpt2_df_derivatives}
%
%For DSRG-MRPT2 the two-electron derivative cannot, in general, be obtained by
%substituting a simple reference pair density into the CASSCF formulas.  The
%second-order energy is a nonvariational functional of the semicanonical
%orbitals, the reference cumulants, the Fock matrix, orbital energy
%denominators, renormalized amplitudes, and renormalized two-electron
%integrals.  The useful way to realize the derivative is therefore to view the
%DSRG-MRPT2 energy evaluation as a scalar computation graph and reverse it to
%obtain adjoints of the DF integral blocks.
%
%\subsection{DSRG-MRPT2 Energy Dependencies}
%
%For a fixed flow parameter $s$, write the energy evaluated by the DSRG-MRPT2
%module schematically as
%\begin{equation}
%E_{\mathrm{DSRG}}
%=
%E_{\mathrm{ref}}
%+
%\mathcal{E}^{(2)}
%\left[
%F, V, \epsilon,
%\gamma, \lambda_2, \lambda_3; s
%\right],
%\label{eq:dsrg_energy_functional}
%\end{equation}
%where $E_{\mathrm{ref}}$ is the bare reference energy and
%$\mathcal{E}^{(2)}$ is the perturbative DSRG correction.  In Forte2 the
%reference orbitals are required to be semicanonical before entering the DSRG
%module.  The DSRG implementation then works in the correlated orbital space
%after any DSRG-level frozen-core and frozen-virtual projections have been
%applied.  With core, active, and virtual correlated orbitals denoted by
%$c$, $a$, and $v$, respectively, the current spin-free implementation builds
%orthonormalized DF blocks
%\begin{equation}
%B^A_{pq} = X_{AP} J^P_{pq},
%\qquad
%X^{\mathrm{T}} X = M^{-1},
%\end{equation}
%for the orbital blocks
%\begin{equation}
%cc,\quad ac,\quad vc,\quad aa,\quad va,
%\end{equation}
%and forms the explicit four-index blocks
%\begin{equation}
%vvaa,\quad aacc,\quad avca,\quad avac,\quad vaaa,\quad aaca,\quad aaaa.
%\end{equation}
%The expensive $ccvv$, $cavv$, and $ccav$ contributions are evaluated
%on-the-fly from three-index DF intermediates rather than stored as full
%four-index arrays.
%
%These implementation details are important for performance, but they should
%not change the integral derivative API.  The integral library should still see
%only the final raw AO derivative weights
%\begin{equation}
%W^P_{\mu\nu}
%\quad\mathrm{and}\quad
%W^M_{PQ},
%\end{equation}
%which are then contracted with $(P|\mu\nu)^{[X]}$ and $(P|Q)^{[X]}$.
%All DSRG-specific response and adjoint logic belongs in the method layer.
%
%\subsection{Adjoints of a DF Four-Index Block}
%
%Let the DSRG reverse pass produce an adjoint
%\begin{equation}
%\bar V_{pqrs}
%\equiv
%\frac{\partial \mathcal{L}}{\partial V_{pqrs}}
%\end{equation}
%for a molecular-orbital two-electron block
%\begin{equation}
%V_{pqrs} = J^P_{pq} N_{PQ} J^Q_{rs},
%\qquad
%J^P_{pq}
%=
%C^*_{\mu p} J^P_{\mu\nu} C_{\nu q}.
%\label{eq:dsrg_mo_df_integral}
%\end{equation}
%Here $\mathcal{L}$ denotes the full DSRG gradient Lagrangian, or just the
%fixed-reference DSRG energy if no orbital, reference, or denominator response
%has been included yet.  Define the metric-applied three-center tensor
%\begin{equation}
%Z^P_{pq} \equiv N_{PQ} J^Q_{pq}.
%\end{equation}
%The differential of a scalar contribution
%\begin{equation}
%\dd \mathcal{L}_{V}
%=
%\bar V_{pqrs} \dd V_{pqrs}
%\end{equation}
%gives the raw MO three-center and metric weights
%\begin{align}
%A^P_{pq}
%&\mathrel{+}=
%\bar V_{pqrs} Z^P_{rs}
%+
%\bar V_{rspq} Z^P_{rs},
%\label{eq:dsrg_mo_j_adjoint}
%\\
%A^M_{PQ}
%&\mathrel{+}=
%-
%\bar V_{pqrs} Z^P_{pq} Z^Q_{rs}.
%\label{eq:dsrg_metric_adjoint}
%\end{align}
%The AO three-center derivative weight is then obtained by back-transforming the
%MO adjoint:
%\begin{equation}
%W^P_{\mu\nu}
%\mathrel{+}=
%C_{\mu p}^* A^P_{pq} C_{\nu q}.
%\label{eq:dsrg_j_backtransform}
%\end{equation}
%Equations \cref{eq:dsrg_mo_j_adjoint,eq:dsrg_metric_adjoint} are the general
%adjoint version of the pair-density formulas derived earlier.  If
%$\bar V_{pqrs} = \frac{1}{2}\Pi_{pq,rs}$ and the pair density is symmetric under
%$(pq)\leftrightarrow(rs)$, they reduce to
%\begin{equation}
%W^P_{pq} = \Pi_{pq,rs} Z^P_{rs},
%\qquad
%W^M_{PQ} = -\frac{1}{2}\Pi_{pq,rs} Z^P_{pq} Z^Q_{rs}.
%\end{equation}
%
%For two-component or otherwise complex orbitals, the same equations are used
%with complex MO adjoints.  The nuclear gradient is real, so the implementation
%should accumulate the real differential,
%\begin{equation}
%\dd \mathcal{L}_{V}
%=
%\mathrm{Re}\left[
%\bar V_{pqrs} \dd V_{pqrs}
%\right],
%\end{equation}
%and contract real AO derivative integrals with the real part of the final AO
%weights.  The MO transformation in \cref{eq:dsrg_mo_df_integral} must use the
%conjugate on the bra coefficient.
%
%\subsection{Why Prefer Raw Metric Adjoints over $B$-Tensor Adjoints}
%
%The existing DSRG code uses orthonormalized DF tensors $B^A_{pq}$ because they
%make contractions simple:
%\begin{equation}
%V_{pqrs} = B^A_{pq} B^A_{rs}.
%\end{equation}
%For gradient work, however, differentiating $B^A_{pq}=X_{AP}J^P_{pq}$ directly
%requires the derivative of the inverse square root of the metric,
%\begin{equation}
%X^{[X]} = \left(M^{-1/2}\right)^{[X]},
%\end{equation}
%which leads to a Sylvester or eigensystem-response problem and introduces
%additional numerical choices.  The raw-metric representation avoids this
%entire issue.  Any adjoint of a four-index block should therefore be converted
%to raw $J$ and $M$ weights using
%\cref{eq:dsrg_mo_j_adjoint,eq:dsrg_metric_adjoint} whenever possible.
%
%If an implementation first obtains adjoints of $B$ blocks, two paths are
%available:
%\begin{enumerate}
%\item Rewrite the differentiated DSRG term in terms of the equivalent
%      four-index block adjoint $\bar V$ and use the raw-metric equations above.
%      This is the preferred route because it automatically includes the metric
%      derivative.
%\item Propagate the $B$ adjoint through $B=XJ$.  This gives the three-center
%      part immediately, but the $X$ adjoint must still be converted into a
%      metric adjoint by differentiating the inverse-square-root relation.  This
%      path should be reserved for cases where a four-index or raw-metric
%      expression is impractical.
%\end{enumerate}
%
%\subsection{Reverse Pass Through DSRG-MRPT2 Blocks}
%
%The method layer should accumulate DSRG derivative weights by reverse
%differentiating each algebraic block used in the energy.  A practical
%decomposition is:
%\begin{enumerate}
%\item Run the usual DSRG-MRPT2 forward pass, but retain the intermediates needed
%      by the reverse pass: semicanonical orbital slices, raw or metric-applied
%      DF blocks, Fock blocks, denominators, amplitudes, renormalization
%      factors, and reference cumulants.
%\item For every explicit energy contraction in
%      $\mathcal{E}^{(2)}$, form adjoints for the participating tensors by the
%      transpose contraction.  In code this mirrors each \texttt{einsum} used in
%      the forward energy.
%\item Accumulate adjoints of the renormalized amplitudes $T_1$, $T_2$,
%      renormalized Fock block $\widetilde F$, and renormalized integrals
%      $\widetilde V$.  The DSRG amplitudes are closed-form algebraic functions,
%      so this is a local reverse pass rather than a large amplitude-response
%      solve.
%\item Convert adjoints of every unrenormalized four-index integral block
%      to raw DF derivative weights using
%      \cref{eq:dsrg_mo_j_adjoint,eq:dsrg_metric_adjoint}.  The same conversion
%      applies to the adjoints generated from the $vvaa$, $aacc$, $avca$,
%      $avac$, $vaaa$, $aaca$, and $aaaa$ blocks.
%\item Treat the on-the-fly $ccvv$, $cavv$, and $ccav$ terms by reversing the
%      three-index contractions inside the loop over the batched core index.
%      This produces adjoints of the temporary bare and renormalized
%      three-index intermediates.  Those adjoints should then be accumulated into
%      the underlying raw DF blocks, or equivalently into the associated
%      four-index adjoints before applying the raw-metric equations.
%\item Add the adjoints from the bare reference energy $E_{\mathrm{ref}}$.
%      These are exactly the CASSCF-like one- and two-particle density
%      contributions derived in the earlier sections.
%\end{enumerate}
%
%The important implementation point is that DSRG-MRPT2 does not need a new
%derivative integral kernel.  It needs a method-specific reverse pass that
%converts the perturbative energy into the same $W^P_{\mu\nu}$ and $W^M_{PQ}$
%weights used by SCF and CASSCF gradients.
%
%\subsection{Renormalization and Denominator Response}
%
%The regularized DSRG amplitudes and renormalized integrals depend explicitly on
%orbital-energy denominators.  For a generic denominator
%\begin{equation}
%\Delta_{ij}^{ab} = \epsilon_i + \epsilon_j - \epsilon_a - \epsilon_b,
%\end{equation}
%the amplitude regularizer is
%\begin{equation}
%R_s(\Delta)
%=
%\frac{1-\exp(-s\Delta^2)}{\Delta},
%\label{eq:dsrg_regularizer}
%\end{equation}
%with the Taylor limit used near $\Delta=0$.  The current utilities also use
%\begin{equation}
%Q_s(\Delta) = 1+\exp(-s\Delta^2)
%\end{equation}
%for the renormalized four-index integrals and the product
%\begin{equation}
%G_s(\Delta) = Q_s(\Delta) R_s(\Delta)
%\end{equation}
%for the on-the-fly three-index intermediates.  The reverse pass therefore
%generates denominator adjoints through
%\begin{align}
%\frac{\partial R_s}{\partial \Delta}
%&=
%\frac{2s\Delta^2\exp(-s\Delta^2)
%-
%\left[1-\exp(-s\Delta^2)\right]}{\Delta^2},
%\label{eq:dsrg_regularizer_deriv}
%\\
%\frac{\partial Q_s}{\partial \Delta}
%&=
%-2s\Delta\exp(-s\Delta^2),
%\\
%\frac{\partial G_s}{\partial \Delta}
%&=
%\frac{\partial Q_s}{\partial \Delta} R_s(\Delta)
%+
%Q_s(\Delta)\frac{\partial R_s}{\partial \Delta}.
%\end{align}
%Near $\Delta=0$, \cref{eq:dsrg_regularizer_deriv} should be evaluated with the
%same Taylor expansion policy as the forward regularizer.
%
%The resulting orbital-energy adjoints have the form
%\begin{equation}
%\bar\epsilon_p
%=
%\sum_\Delta \bar\Delta\,
%\frac{\partial \Delta}{\partial \epsilon_p}.
%\end{equation}
%They are not derivative-integral weights by themselves.  They must be folded
%into the orbital and semicanonical response equations, or into an equivalent
%Lagrangian enforcing the definition of the orbital energies.  Omitting these
%terms gives a fixed-denominator derivative, not a full DSRG-MRPT2 gradient.
%
%\subsection{Semicanonical Orbital Response}
%
%DSRG-MRPT2 in Forte2 assumes that the parent method has delivered
%semicanonical orbitals.  The active orbitals and active cumulants are also
%rotated by the semicanonical transformation before the perturbative energy is
%evaluated.  Thus the full derivative contains response of
%\begin{equation}
%C_{\mathrm{semi}},
%\qquad
%\epsilon_p,
%\qquad
%\gamma,
%\qquad
%\lambda_2,
%\qquad
%\lambda_3.
%\end{equation}
%There are two reasonable formulations:
%\begin{enumerate}
%\item Add semicanonicalization constraints to the DSRG Lagrangian.  The
%      constraints enforce zero off-diagonal Fock blocks inside the
%      semicanonical subspaces and define the orbital energies.  The
%      corresponding adjoint variables absorb the derivatives of
%      $C_{\mathrm{semi}}$, $U_{\mathrm{act}}$, and $\epsilon$.
%\item Reformulate as much of the DSRG energy as possible in a rotationally
%      covariant block form, and only introduce denominators through explicitly
%      constrained Fock eigenvalues.  This reduces, but does not eliminate, the
%      need for a denominator-response adjoint.
%\end{enumerate}
%The first route is the most direct match to the current code because the
%forward calculation is already semicanonical and denominator based.
%
%\subsection{Reference and CASSCF Response}
%
%For a CASSCF parent, DSRG-MRPT2 introduces a target energy
%\begin{equation}
%E_T = E_{\mathrm{ref}} + \mathcal{E}^{(2)}
%\end{equation}
%that is not stationary with respect to the parent CASSCF orbital and CI
%variables.  The CASSCF Z-vector equations in
%\cref{sec:casscf_zvector} should therefore be used with a DSRG-specific
%right-hand side:
%\begin{equation}
%A_{mn}^{\mathrm{CAS}} z_n^{\mathrm{CAS}}
%=
%-
%\frac{\partial E_T}{\partial y_m^{\mathrm{CAS}}}.
%\label{eq:dsrg_cas_zvector_rhs}
%\end{equation}
%The right-hand side is assembled from the DSRG reverse pass:
%\begin{equation}
%\frac{\partial E_T}{\partial y_m^{\mathrm{CAS}}}
%=
%\bar F_{pq}\frac{\partial F_{pq}}{\partial y_m^{\mathrm{CAS}}}
%+
%\bar V_{pqrs}\frac{\partial V_{pqrs}}{\partial y_m^{\mathrm{CAS}}}
%+
%\bar\gamma_{uv}\frac{\partial \gamma_{uv}}{\partial y_m^{\mathrm{CAS}}}
%+
%\bar\lambda_{uvxy}\frac{\partial \lambda_{uvxy}}{\partial y_m^{\mathrm{CAS}}}
%+
%\bar\lambda_{uvwxyz}
%\frac{\partial \lambda_{uvwxyz}}{\partial y_m^{\mathrm{CAS}}}
%+
%\bar\epsilon_p
%\frac{\partial \epsilon_p}{\partial y_m^{\mathrm{CAS}}}.
%\end{equation}
%After solving the CASSCF response problem, the DSRG relaxed densities and
%adjoint DF weights are added to the bare CASSCF reference-density contribution.
%
%If DSRG reference relaxation is used, the differentiated objective changes.
%For a one-shot or iterative relaxed-reference calculation, the Lagrangian must
%also enforce the relaxed-reference eigenvalue or residual equations generated
%by the effective Hamiltonian $\bar H$.  Equivalently, after convergence, one can
%treat the relaxed reference as the stationary state of the final effective
%Hamiltonian, but then the adjoints of the $\bar H$ one- and two-body
%intermediates must be included in the same reverse pass.
%
%\subsection{Frozen Core and Frozen Virtual Orbitals}
%
%The DSRG module may apply additional frozen-core or frozen-virtual projections
%relative to the parent method.  In the current layout, the perturbative
%correlation part is built only from the remaining correlated core, active, and
%virtual slices.  Therefore:
%\begin{enumerate}
%\item Frozen orbitals excluded from the DSRG correlated space should not appear
%      in DSRG amplitude denominators or perturbative DF blocks.
%\item The bare reference energy still contains the appropriate frozen-core
%      contribution inherited from the parent reference.  Its explicit
%      derivative is part of the reference-density gradient, not the DSRG
%      correlation adjoint.
%\item If the frozen core is parent-relaxed, the frozen-core parent response
%      described in \cref{eq:frozen_core_zvector,eq:frozen_core_gradient} must be
%      included in the total Z-vector system.  The DSRG perturbative right-hand
%      side contributes to the parent orbital-response equations whenever the
%      DSRG energy depends on rotations involving the frozen parent orbitals.
%\item Frozen virtual orbitals require the analogous projection bookkeeping.
%      They are absent from the DSRG amplitudes, but parent orbital response may
%      still be required if their definition changes with geometry.
%\end{enumerate}
%
%\subsection{Suggested Forte2 Realization}
%
%A staged implementation can be organized as follows:
%\begin{enumerate}
%\item Extend the DF builder or add a DSRG-specific helper that can provide raw
%      MO three-center blocks $J^P_{pq}$ and the inverse metric $N_{PQ}$, in
%      addition to the current orthonormalized $B^A_{pq}$ blocks.  This avoids
%      differentiating $M^{-1/2}$.
%\item Refactor the DSRG-MRPT2 energy routines so that each contraction has a
%      matching reverse routine.  The reverse routines should accumulate
%      adjoints to $F$, $V$, $J$, $\epsilon$, and the reference cumulants.
%\item Implement reverse versions of the DSRG regularization utilities
%      \texttt{compute\_T1\_block}, \texttt{compute\_T2\_block},
%      \texttt{renormalize\_V\_block}, and \texttt{renormalize\_3index}.  These
%      routines should return both tensor adjoints and orbital-energy adjoints.
%\item For the on-the-fly $ccvv$, $cavv$, and $ccav$ paths, accumulate local
%      adjoints inside the existing loop over the batched core index.  This
%      keeps memory comparable to the forward code and avoids materializing the
%      full four-index blocks.
%\item Back-transform the final MO three-center adjoints to AO weights and add
%      the metric adjoint to $W^M_{PQ}$.
%\item Solve the required orbital, semicanonical, CASSCF, and frozen-core
%      Z-vector equations.  The relaxed one-electron, two-electron, overlap, and
%      DF derivative weights are then passed to the same integral derivative
%      kernels used by the lower-level methods.
%\end{enumerate}
%
%\subsection{Validation Strategy}
%
%The DF derivative should be validated in layers:
%\begin{enumerate}
%\item First test only the explicit DF integral derivative by holding orbitals,
%      denominators, amplitudes, and cumulants fixed and finite-differencing the
%      DSRG scalar energy with respect to displaced three-center and metric
%      integrals.
%\item Next compare fixed-reference molecular gradients against finite
%      differences of the DSRG-MRPT2 energy evaluated at displaced geometries,
%      including denominator and semicanonical response but without reference
%      relaxation.
%\item Then enable CASSCF response and state-averaged parent references.  The
%      Z-vector right-hand side should be tested against finite differences of
%      individual-state and state-averaged DSRG energies.
%\item Finally test DSRG reference relaxation, frozen-core variants, and the
%      two-component complex path.  For complex calculations, the imaginary part
%      of the final gradient should be numerically negligible and the reported
%      gradient should be the real differential.
%\end{enumerate}
%
%\section{Implementation Decomposition}
%\label{sec:implementation_decomposition}
%
%The implementation should separate method-independent derivative integral
%kernels from method-specific density and response assembly.  The integral
%library should know how to contract derivative integral buffers with supplied
%weights.  It should not know whether those weights came from RHF, UHF, GHF,
%CASSCF, state-averaged CASSCF, DSRG-MRPT2, or a Z-vector Lagrangian.
%
%\subsection{Integral-Library Functions}
%
%The following functions are method independent and belong in the integral
%library.
%
%\paragraph{Center validation.}
%The one-electron derivative code already validates that the basis centers and
%the molecular charge centers are consistent.  The same logic should be reused
%or generalized for three-center and two-center DF derivative kernels:
%\begin{equation}
%\texttt{validate\_deriv\_centers(basis, atoms, basis\_name)}.
%\end{equation}
%For a three-center kernel, both the auxiliary basis and the orbital basis sets
%must have centers matching the nuclear coordinate list.  This is especially
%important when the auxiliary basis has different shell content but the same
%atomic centers as the orbital basis.
%
%\paragraph{Three-center derivative contraction.}
%Implement a function with the logical signature
%\begin{equation}
%\texttt{coulomb\_3c\_deriv(aux\_basis, basis1, basis2, W3, atoms)}
%\rightarrow
%\mathbf{G}^{(3c)}.
%\end{equation}
%It computes
%\begin{equation}
%G^{(3c)}_{A\alpha}
%=
%W^P_{\mu\nu}
%(P|\mu\nu)^{[A\alpha]}.
%\label{eq:coulomb_3c_deriv_contract}
%\end{equation}
%The expected weight shape is
%\begin{equation}
%W3[P,\mu,\nu],
%\end{equation}
%matching the raw three-center integral ordering $(P|\mu\nu)$.  The function
%should:
%\begin{enumerate}
%\item check that \texttt{W3.shape == (naux, nbasis1, nbasis2)};
%\item initialize a first-derivative Libint2 Coulomb engine with
%      \texttt{BraKet::xs\_xx};
%\item loop over auxiliary and orbital shell triplets;
%\item compute derivative buffers for the shell centers of $P$, $\mu$, and
%      $\nu$;
%\item map each derivative buffer component to the atom on which that shell is
%      centered;
%\item contract each buffer element with the corresponding weight and accumulate
%      into a real gradient vector of length $3N_{\mathrm{atom}}$.
%\end{enumerate}
%For complex methods such as complex GHF, either provide a complex-weight
%overload or accept a complex ndarray and accumulate
%\begin{equation}
%G^{(3c)}_{A\alpha}
%=
%\mathrm{Re}
%\left[
%W^P_{\mu\nu}
%(P|\mu\nu)^{[A\alpha]}
%\right].
%\end{equation}
%The derivative integrals are real for a real scalar AO and auxiliary basis; the
%real part enforces the real-valued energy-gradient convention.
%
%\paragraph{Two-center metric derivative contraction.}
%Implement a function with the logical signature
%\begin{equation}
%\texttt{coulomb\_2c\_deriv(aux\_basis1, aux\_basis2, W2, atoms)}
%\rightarrow
%\mathbf{G}^{(2c)}.
%\end{equation}
%It computes
%\begin{equation}
%G^{(2c)}_{A\alpha}
%=
%W^M_{PQ}(P|Q)^{[A\alpha]}.
%\label{eq:coulomb_2c_deriv_contract}
%\end{equation}
%The expected weight shape is
%\begin{equation}
%W2[P,Q].
%\end{equation}
%The function should:
%\begin{enumerate}
%\item check that \texttt{W2.shape == (naux1, naux2)};
%\item initialize a first-derivative Libint2 Coulomb engine with
%      \texttt{BraKet::xs\_xs};
%\item loop over auxiliary shell pairs;
%\item compute derivative buffers for the two auxiliary shell centers;
%\item map the derivative components to atoms and contract with $W^M_{PQ}$;
%\item return a real gradient vector of length $3N_{\mathrm{atom}}$.
%\end{enumerate}
%If the two auxiliary bases are the same and $W^M$ is symmetric, the
%implementation may use a symmetry-reduced shell-pair loop with off-diagonal
%factors.  A full shell-pair loop is simpler and is the safer first
%implementation.
%
%\paragraph{Combined DF derivative contraction.}
%Add a convenience wrapper
%\begin{equation}
%\texttt{df\_coulomb\_deriv(aux\_basis, basis, W3, W2, atoms)}
%\rightarrow
%\mathbf{G}^{(2)}.
%\end{equation}
%It should return
%\begin{equation}
%\mathbf{G}^{(2)}
%=
%\mathbf{G}^{(3c)}
%+
%\mathbf{G}^{(2c)}.
%\end{equation}
%This wrapper should not build $W3$ or $W2$; it only dispatches the two
%derivative contractions.  Keeping the wrapper weight-driven makes it reusable
%for all wave-function models.
%
%\paragraph{Optional by-shell variants.}
%For large systems, provide block variants analogous to the existing
%three-center integral slices:
%\begin{align}
%&\texttt{coulomb\_3c\_deriv\_by\_shell(aux\_basis, basis1, basis2, shell\_slices, W3, atoms)},\\
%&\texttt{coulomb\_2c\_deriv\_by\_shell(aux\_basis1, aux\_basis2, shell\_slices, W2, atoms)}.
%\end{align}
%These are useful when $W3$ is generated in auxiliary or AO blocks and the full
%$N_{\aux}N_{\ao}^2$ tensor cannot be stored.  The by-shell functions should
%accumulate into a caller-provided gradient vector or return the block
%contribution for the caller to sum.
%
%\paragraph{Derivative integral tensor builders for testing.}
%The production path should contract derivative buffers immediately, but tests
%are easier if small derivative tensors can be materialized:
%\begin{align}
%&\texttt{coulomb\_3c\_deriv\_tensor(aux\_basis, basis1, basis2, atoms)}
%\rightarrow
%T_{A\alpha,P\mu\nu},\\
%&\texttt{coulomb\_2c\_deriv\_tensor(aux\_basis1, aux\_basis2, atoms)}
%\rightarrow
%U_{A\alpha,PQ}.
%\end{align}
%These functions should be optional, test-only, or clearly documented as
%small-system utilities because their storage scales poorly.
%
%\paragraph{Nanobind exposure and stubs.}
%The public C++ functions should be exposed in the integrals API, with
%overloads for identical basis arguments where useful.  After binding changes,
%the Python stubs should be regenerated so that the low-level
%\texttt{forte2.ints} functions are visible to Python-side method code and
%tests.
%
%\subsection{Method-Specific Functions}
%
%The following functions should live outside the integral library, in SCF,
%CASSCF, DSRG, or shared gradient modules.
%
%\paragraph{Metric application and DF intermediates.}
%Build or apply the inverse Coulomb metric:
%\begin{equation}
%Z^P_{\mu\nu}=N_{PQ}J^Q_{\mu\nu}.
%\end{equation}
%This is method independent mathematically, but it belongs near the DF/Fock
%builder because it depends on the chosen metric treatment: Cholesky solves,
%eigenvalue truncation, or an existing stored inverse-square-root representation.
%The derivative integral library should receive only the final weights.
%
%\paragraph{SCF weight builders.}
%For each SCF reference, construct $W3$ and $W2$ from the converged density:
%\begin{align}
%\texttt{build\_rhf\_df\_gradient\_weights} &:
%P=2D,\quad
%\Pi=P P-\frac{1}{2}P P_{\mathrm{exchange}},\\
%\texttt{build\_uhf\_df\_gradient\_weights} &:
%D^\alpha,\ D^\beta,\quad
%\Pi=P P-D^\alpha D^\alpha_{\mathrm{exchange}}
%-D^\beta D^\beta_{\mathrm{exchange}},\\
%\texttt{build\_rohf\_df\_gradient\_weights} &:
%\text{same pair density as UHF, but with common spatial orbitals},\\
%\texttt{build\_ghf\_df\_gradient\_weights} &:
%\gamma^{\alpha\alpha},\gamma^{\alpha\beta},
%\gamma^{\beta\alpha},\gamma^{\beta\beta}.
%\end{align}
%These functions produce the method-independent tensors
%\begin{equation}
%W3[P,\mu,\nu],
%\qquad
%W2[P,Q],
%\end{equation}
%then call the integral-library derivative contractions.
%
%\paragraph{One-electron and overlap-density builders.}
%A full gradient driver also needs method-specific builders for:
%\begin{enumerate}
%\item the one-particle density that contracts with $h_{\mu\nu}^{[X]}$;
%\item the energy-weighted density or orbital Lagrange multiplier that contracts
%      with $S_{\mu\nu}^{[X]}$;
%\item the nuclear repulsion derivative.
%\end{enumerate}
%These are not part of the two-electron DF derivative library.  They should be
%implemented in the method gradient layer so that ROHF canonicalization,
%complex GHF conjugation, and CASSCF response conventions are handled in one
%place.
%
%\paragraph{CASSCF and SA-CASSCF density builders.}
%For CASSCF, construct the internal-space pair density from core occupancy and
%active-space RDMs, or directly accumulate the equivalent DF weights from the
%core-core, core-active, and active-active energy pieces:
%\begin{equation}
%\Pi_{pq,rs}
%=
%\left\langle
%\hat{E}_{pq}\hat{E}_{rs}
%-
%\delta_{qr}\hat{E}_{ps}
%\right\rangle.
%\end{equation}
%For state-averaged CASSCF, build
%\begin{equation}
%\Pi^\sa_{pq,rs}=\sum_I w_I\Pi^I_{pq,rs}.
%\end{equation}
%The method code then transforms the MO-basis three-center weight back to the
%AO basis:
%\begin{equation}
%W^P_{\mu\nu}=C^*_{\mu p}W^P_{pq}C_{\nu q}.
%\end{equation}
%
%\paragraph{Z-vector and relaxed-density builders.}
%The CASSCF gradient layer must solve the Z-vector equations described in
%\cref{sec:casscf_zvector} when the target is not stationary in the optimized
%CASSCF variables.  The outputs are relaxed one- and two-particle densities:
%\begin{equation}
%\overline{\gamma}=\gamma^T+\gamma^Z,
%\qquad
%\overline{\Pi}=\Pi^T+\Pi^Z.
%\end{equation}
%Those relaxed densities are then passed to the same CASSCF DF weight builder.
%Frozen-core response, if parent-relaxed rather than constrained, also belongs
%in this method layer.
%
%\paragraph{DSRG-MRPT2 adjoint builders.}
%For DSRG-MRPT2, construct $W3$ and $W2$ from the adjoint reverse pass described
%in \cref{sec:dsrg_mrpt2_df_derivatives}.  The method-specific code should:
%\begin{enumerate}
%\item retain the raw or metric-applied MO three-center blocks needed to avoid
%      differentiating the orthonormalized $B$ tensor;
%\item reverse the explicit and on-the-fly DSRG energy contractions to obtain
%      adjoints of $F$, $V$, $J$, $\epsilon$, and the reference cumulants;
%\item include reverse derivatives of the DSRG regularizers so orbital-energy
%      denominator response enters the orbital and semicanonical Z-vector
%      equations;
%\item add CASSCF, state-averaged CASSCF, relaxed-reference, and frozen-core
%      response corrections when those variants define the target energy;
%\item back-transform the final three-center adjoints to AO weights and pass
%      them, together with the metric adjoint, to the integral-library kernels.
%\end{enumerate}
%
%\paragraph{Gradient drivers.}
%Finally, each method should provide a high-level driver that assembles:
%\begin{equation}
%\mathbf{G}
%=
%\mathbf{G}_{\mathrm{nuc}}
%+
%\mathbf{G}_{1e}
%+
%\mathbf{G}_{2e,\df}
%+
%\mathbf{G}_{S}.
%\end{equation}
%Here $\mathbf{G}_{2e,\df}$ is supplied by the integral-library derivative
%contractions once the method layer has generated $W3$ and $W2$.
%
%\section{Derivative Integral Kernel Mapping}
%
%For each nuclear coordinate $X=R_{A\alpha}$, the implementation target is
%\begin{equation}
%G^{(2)}_{A\alpha}
%=
%W^P_{\mu\nu}(P|\mu\nu)^{[A\alpha]}
%+
%W^M_{PQ}(P|Q)^{[A\alpha]}.
%\end{equation}
%The three-center derivative integral kernel should accumulate derivatives with
%respect to all centers in the shell triplet $(P|\mu\nu)$:
%\begin{equation}
%(P|\mu\nu)^{[A\alpha]}
%=
%\frac{\partial}{\partial R_{A\alpha}}
%(P|\mu\nu).
%\end{equation}
%If the auxiliary shell $P$ is centered on atom $A$, its derivative buffer
%contributes to $G^{(2)}_{A\alpha}$.  If the orbital shells $\mu$ or $\nu$ are
%centered on atom $A$, their derivative buffers also contribute to
%$G^{(2)}_{A\alpha}$.
%
%Likewise, the two-center metric derivative kernel should accumulate derivatives
%with respect to both centers in $(P|Q)$:
%\begin{equation}
%(P|Q)^{[A\alpha]}
%=
%\frac{\partial}{\partial R_{A\alpha}}
%(P|Q).
%\end{equation}
%Because $W^M_{PQ}$ is symmetric, the metric contribution may use the full
%$P,Q$ shell loop directly or a symmetry-reduced shell loop with the usual
%off-diagonal factor.
%
%\section{Summary of the Working Equations}
%
%The density-fitted two-electron gradient contribution can be implemented with
%the following sequence:
%\begin{enumerate}
%\item Build or apply the inverse Coulomb metric
%      $N_{PQ}=(P|Q)^{-1}$.
%\item Build the three-center integrals in the basis in which the pair density
%      is available:
%      \begin{equation}
%      J^P_{pq}=C^*_{\mu p}(P|\mu\nu)C_{\nu q}.
%      \end{equation}
%\item Apply the inverse metric:
%      \begin{equation}
%      Z^P_{pq}=N_{PQ}J^Q_{pq}.
%      \end{equation}
%\item Build the derivative weights:
%      \begin{align}
%      W^P_{pq}
%      &=
%      \Pi_{pq,rs}Z^P_{rs},
%      \\
%      W^M_{PQ}
%      &=
%      -\frac{1}{2}
%      \Pi_{pq,rs}Z^P_{pq}Z^Q_{rs}.
%      \end{align}
%      For nonvariational methods such as DSRG-MRPT2, replace this
%      pair-density step by the adjoint equations
%      \cref{eq:dsrg_mo_j_adjoint,eq:dsrg_metric_adjoint}.
%\item Transform the three-center weight to the AO basis:
%      \begin{equation}
%      W^P_{\mu\nu}=C^*_{\mu p}W^P_{pq}C_{\nu q}.
%      \end{equation}
%\item Contract first-derivative integral buffers:
%      \begin{equation}
%      E_2^{[X]}
%      =
%      W^P_{\mu\nu}(P|\mu\nu)^{[X]}
%      +
%      W^M_{PQ}(P|Q)^{[X]}.
%      \end{equation}
%\end{enumerate}

\end{document}
